% Build from docs with pdflatex type-system.tex until the references settle.
\documentclass[11pt]{article}
\usepackage[T1]{fontenc}
\usepackage{lmodern}
\usepackage[margin=1.05in]{geometry}
\usepackage{amsmath,amssymb,amsthm}
\usepackage[hyphens]{url}
\usepackage[hidelinks]{hyperref}
\hypersetup{pdftitle={The type system of Coalton}}
\setcounter{tocdepth}{2}
\setlength{\emergencystretch}{1.5em}
\raggedbottom

\newtheorem{theorem}{Theorem}[section]
\newtheorem{lemma}[theorem]{Lemma}
\newtheorem{proposition}[theorem]{Proposition}
\theoremstyle{definition}
\newtheorem{example}[theorem]{Example}
\newtheorem{invariant}[theorem]{Invariant}
\newcommand{\code}[1]{\texttt{#1}}
\newcommand{\src}[1]{\nolinkurl{#1}}
\newcommand{\fv}{\operatorname{fv}}
\newcommand{\dom}{\operatorname{dom}}
\newcommand{\pack}{\operatorname{out}}
\newcommand{\Fn}{\mathsf{Fn}}
\newcommand{\Pack}{\mathsf{Pack}}
\newcommand{\Bool}{\mathsf{Boolean}}
\newcommand{\Int}{\mathsf{Integer}}
\newcommand{\Str}{\mathsf{String}}
\newcommand{\Unit}{\mathsf{Unit}}
\newcommand{\Void}{\mathsf{Void}}
\newcommand{\List}{\mathsf{List}}
\newcommand{\Dict}{\mathsf{Dict}}
\newcommand{\sel}{\operatorname{select}}
\newcommand{\clo}{\operatorname{cl}}
\newcommand{\eqs}{\doteq}
\newcommand{\mgu}{\operatorname{mgu}}
\newcommand{\oneway}{\mathrel{\triangleright}}
\newcommand{\given}{\,;\,}
\newcommand{\chk}{\mathrel{\Leftarrow}}
\newcommand{\arr}{\mathbin{\to}}
\newcommand{\rrule}[3]{\frac{\displaystyle #2}{\displaystyle #3}
                       \;{\scriptstyle[\textsf{#1}]}}
% Keep identifier-heavy source notes in a short, ragged-right paragraph.
\long\def\sources#1\par{\par\smallskip\begingroup\small\raggedright
                       \noindent\textit{Source:} #1\par\endgroup}

\title{The type system of Coalton}
\author{}
\date{Source revision \texttt{8f69e30780f4}, based on \texttt{34e112ad}}
\begin{document}
\maketitle

{\centering
\fbox{An automated AI system created this document and updates it periodically.}\par}

\begin{abstract}
Coalton checks expressions against expected types while inferring substitutions,
class constraints and field selections. Three operations that look like solving
the same equation have different authority: unification may determine a caller's
unknowns, instance matching must not, and a constructor test may refine an
index only on the path on which the test succeeds. We specify these operations,
where inference calls them, and what leaves each scope. We prove first-order
unification, fixed-environment evidence and conditional improvement results,
and isolate a sharing invariant at generalization. We prove progress and
preservation for a pure, nonrecursive indexed core with exhaustive cases.
The connection to the checker covers a monomorphic fragment and a family of
explicitly quantified eliminators, not arbitrary GADT inference or generated
Common Lisp. A compiled instance retains its evidence interface across
extensions. We show why admission justifies agreement on shared inputs
under its prerequisite hypotheses and how recursive instance search
can reach a base case even when its arguments grow, without treating
that growth as evidence.
\end{abstract}
\tableofcontents

\section{What is being specified}

The input is a parsed, macroexpanded (where the parser expands macros),
lexically renamed Coalton program, in a fixed compiler environment. Symbols
include their packages; different fresh names with the same printed spelling
are different variables. Source examples use the usual Coalton syntax;
formulas use an abstract syntax in which, for instance, $D\,s\,t$ is
left-associated type application, not a call. A bar denotes an ordered finite
list, possibly empty; $P\mathbin{+\!\!+}Q$ concatenates two lists in order.
Substitutions act pointwise on lists. Finite sets are used
when order and multiplicity are immaterial. $\fv$ means free \emph{type}
variables, including free skolems, not term variables. Spelling, diagnostic
locations and alias histories do not justify type equations.

There are three levels of statement here:
\begin{itemize}
\item A \emph{meaning} rule says what a type or piece of evidence licenses.
It need not say how to find it. In particular, an equality given by a pattern
is not a unification request from its body.
\item A \emph{checking} rule or procedure describes a successful compiler path;
failed premises reject that path unless a fallback is specified. We specify
type-relevant decisions, not the choice of diagnostic or the layout of every
typed node. Inference is not defined as ``find any declarative derivation.''
\item A \emph{theorem} identifies its domain and hypotheses. In particular,
only Section~\ref{sec:safety} assigns an operational semantics to terms.
No theorem silently assumes the whole compiler implements that semantics.
\end{itemize}
The global signature $\Sigma$ contains nominal constructors and their kinds,
aliases, data constructors and their value schemes, structs and fields,
classes, instances, and global value schemes. The compiler resolves lexical
control targets, checks type/struct/alias groups, then classes, installs
instance heads (including generated and derived instances), finalizes their
dictionary signatures as a batch, checks value definitions, checks instance
bodies, and processes explicit specializations.
Analysis of the typed unit, dictionary translation, optimizations and host
compilation follow. Thus instance selection is relative to a signature, not
to the textual prefix ending at a call. Source anchors below are relative to
\src{src/typechecker/} unless a path is given.

\sources \src{src/entry.lisp}: \code{\%entry-point}, \src{entry-point};
\src{tc-env.lisp}: \src{tc-env-lookup-value}, \src{tc-env-parser-env};
\src{src/parser/renamer.lisp}: \src{rename-variables}.

\section{Type language and its sorts}\label{sec:types}

The following grammar describes internal objects; position and kind checks
restrict their combinations. $D$ is a nominal type constant, $C$ a class.
$a$ is a flexible type variable, $s$ a rigid variable (skolem), and $g_i$ a
generic index in a stored scheme. Write $\alpha$ for any type variable when
rigidity is unimportant.
\[
\begin{aligned}
 \kappa &::= * \mid \kappa\arr\kappa \mid u
                       && u\text{ a kind unknown}\\
 t &::= a \mid s \mid g_i \mid D \mid t\,t
        \mid \Fn(\bar t;K;o;\bar t)
        \mid \Pack(\bar t) \mid \forall\bar\alpha.\,Q\Rightarrow t\\
 K &::= \{\,{:}k_i:t_i\,\}_{i=1}^m
          && o\in\{\mathsf{closed},\mathsf{open}\}\\
 p &::= C\,\bar t,\qquad Q ::= [p_1,\dots,p_n].
\end{aligned}
\]
In $\Fn(\bar A;K;o;\bar B)$, $\bar A$ is the positional input vector,
$K$ a finite map of optional keyword inputs (stored in keyword-name order),
and $\bar B$ the output vector. An open inferred function can also carry an
\emph{origin}, a variable used to propagate new keywords; it is not a
program-visible argument or an additional type parameter of the function.
Write $\bar A\arr t$ for a closed keyword-free function with result $t$.
One $A * B \arr C$ call consumes two positional arguments;
$A\arr(B\arr C)$ consumes one and returns a function. A nullary function
is not its result. $\mathsf{Arrow}:*\arr *\arr *$ permits partially applied
arrows at higher kind. A saturated \code{Arrow} normalizes to a unary
function; comparison against an application spine can re-express a closed,
keyword-free, one-input/one-output function as an arrow application if neither
endpoint is a \emph{bare} nested $\forall$.

\paragraph{Kinding.}
For $\Delta$ a finite assignment of kinds to variables, the leaves have their
declared kinds ($g_i$ uses its scheme's $i$th kind). The central rules are
\[
 \rrule{K-app}{\Delta\vdash f:\kappa\arr\lambda\quad
                   \Delta\vdash a:\kappa}{\Delta\vdash f\,a:\lambda}
 \qquad
 \rrule{K-fn}{\Delta\vdash \bar A,\operatorname{ran}K,\bar B:*}
             {\Delta\vdash\Fn(\bar A;K;o;\bar B):*}.
\]
A pack of $*$ components has kind $*$; a nested $\forall$ has kind $*$ when
its body has kind $*$ and its predicates have the class parameter kinds
under the extended $\Delta$. A class is a relation on its parameter types,
not a type of kind $*$. Predicate arity must equal the class arity. Kinding
collects equations (give an unknown head the kind
$\kappa_{\rm arg}\arr\kappa_{\rm expected}$), solves by first-order
unification on $*,\arr,u$ with an occurs check, and replaces still
undetermined kind variables by $*$ at the appropriate declaration/group
boundary. There is higher-kinded inference, not polymorphism over kinds in
stored schemes. Mutually dependent data declarations share preliminary
constructor kinds; aliases cannot form a cycle. Aliases expand by application
of their parameters, retaining a printing history. The post-expansion type,
not an alias's name, is checked.

\paragraph{Results are not tuples.}
Define the result of an output vector by
\[
 \pack([])=\Pack([])=\Void,\qquad
 \pack([t])=t,\qquad
 \pack([t_1,\ldots,t_n])=\Pack([t_1,\ldots,t_n])\quad(n\ge2).
\]
Writing $\pack^{-1}(t)$ means the component list of a pack and $[t]$
otherwise. A singleton output that becomes a pack under substitution is
spliced into its function's outputs. A \code{Tuple} is one nominal value;
$\Unit$ has a value, whereas $\Void$ here means \emph{zero results}, not
an uninhabited value type. A plain variable $a^v$ ranges over single-value
types. A result-capable variable $a^r$ may also stand for zero or several
outputs. Source \code{Values} on a binder sets that flag at kind $*$; it
is not a second base kind. A scoped such binder may occupy a whole function
or \code{lisp} output slot, not an input, class argument, type argument or
element of a multi-output list. Thus one signature can forward a call's
entire output without giving \code{List} a type of lists of packs.

\paragraph{Two representations of polymorphism.}
A stored value scheme $\sigma=\forall\bar a.\,Q\Rightarrow t$ is a kind
vector and a qualified body with selected variables replaced by $g_i$.
$\operatorname{inst}(\sigma)$ replaces all occurrences of each $g_i$ by
\emph{one} fresh flexible variable of its kind and result capability, shared
across predicates and body. Free variables, e.g.\ a weak variable in a local
scheme, are \emph{not} freshened. Quantification picks eligible nonrigid
variables actually occurring in the qualified type. Explicit binder order
and source names are preserved for scope/printing; generalized open origins
are refreshed separately, consistently within an instantiation. A nested
$\forall$ instead binds named-by-identity variables in its own qualified body:
substitution excludes their domains, and instantiation freshens binders to
avoid capture. We always choose fresh identities when entering a binder.

Ordinary annotations implicitly bind their unscoped free variables where
schemes are formed. Explicit \code{forall} introduces a lexical scope;
unlisted free variables must already be in scope. E.g.\ an explicit
\code{:a} in a function declaration is available in a \code{the} or
\code{lisp} inside the definition; a printed \code{:a} in an implicitly
quantified declaration does not establish that lexical name. Class/instance
head variables and explicitly scoped method binders also have body scopes.
A nested predicate context requires an explicit \code{forall}. A bare
quantified type can be a function input, single output or nominal field,
but not a type application operand, a class argument or a multi-value
component; a nested $\forall$ cannot quantify a pack directly. These are
positional checks (also performed on expanded types), not a ban on
functions that contain polymorphic slots. A root alias expanding to
$\forall$ has its freshened binders and predicates lifted into the scheme.

\sources \src{types.lisp}: \src{apply-type-argument}, \src{output-types-result-type},
\src{normalize-function-output-types}; \src{kinds.lisp}: \src{kmgu},
\src{kind-monomorphize-subs}; \src{parse-type.lisp}: \src{infer-type-kinds},
\src{check-result-variable-positions}, \src{parse-qualified-type-info},
\src{check-no-impredicative-forall-types}; \src{scheme.lisp}: \src{quantify},
\src{fresh-inst}; \src{predicate.lisp}: \src{instantiate-forall-ty}.

\section{Equality, unification and matching}\label{sec:unification}

Write $St$ for simultaneous substitution: a variable lookup returns the
first right-hand side for that variable, not a recursive application of
the substitution to itself. On ordinary compound types substitution is
structural. Composition $T\circ S$ maps each $a\mapsto t$ in $S$ to
$a\mapsto Tt$, followed by $T$'s entries. Thus on ordinary trees
$(T\circ S)t=T(St)$. Function outputs and arrow applications normalize;
if substitution maps an open function's origin to another function, the
updated function replaces the old one. These representation operations are
why it is useful to state the unifier theorem on a specified grammar.

$t=u$ in the algorithm below denotes \src{ty=}: the same constructor
symbol and kind, or the same variable ID, kind, rigidity and result flag,
or structural equality. Aliases, diagnostic existential flags and origin
identity do not enter this test. For functions it compares full arities,
keyword entries, openness and outputs; for nested binders its fast path
compares kinds/flags in order and bodies opened with shared fresh indices.
Predicate equality compares the class and all arguments. The order of
predicates matters to \src{qualified-ty=}; the separate binding scheme
comparison \src{ty-scheme=} instead seeks a kind- and flag-respecting
renaming on occurrences and compares predicate \emph{bags}, tracking binder
scopes. Equality is not class entailment or observational equivalence.

\paragraph{Ordinary equations.}
An equation $t\eqs u$ extends $S$ by
$\mgu(St,Su)\circ S$. Equal constants need no binding; unequal constant
heads fail. Applications decompose head then argument, using the head
solution before comparing arguments. In each binding of $a$ to $u$:
\begin{enumerate}
 \item Identical variables succeed unchanged. A rigid variable is never
 bound; a rigid/flexible pair is oriented towards the rigid variable if
 admissible, otherwise it fails.
 \item Kinds must agree. A flexible variable cannot bind to a bare $\forall$.
 Only a result-capable variable can bind to a pack. Between flexible variables
 of different result capability the ordinary variable survives; an ordinary
 variable cannot bind to a result-capable skolem.
 \item If $a$ occurs free in the proposed nonvariable right-hand side, fail;
 otherwise emit $[a\mapsto u]$.
\end{enumerate}
Packs unify componentwise at equal lengths. Functions require equal input
arity, keyword sets and openness; unify positional slots, corresponding
keyword slots, and $\pack(\bar B_1)\eqs\pack(\bar B_2)$ in sequence.
Unequal output counts are tried only if one side is a singleton
result-capable variable. If both functions carry origins, also unify the
origins. Keyword \emph{projection} at a use site is a separate operation
(Section~\ref{sec:terms}). For two $\forall$ types, after the equality fast
path, check equal binder and predicate counts. Open one side rigidly and the
other flexibly; unify corresponding predicate arguments (same class in each
position) and then the bodies. The fresh flexibles must map bijectively to
the skolems with matching kinds/flags, and no remaining substitution may
contain these skolems. Drop the temporary domains. This permits renaming
and unification of free variables, not arbitrary subsumption under arrows.

\paragraph{One-way equations.}
$t\oneway u\leadsto M$ means that $t$ is the \emph{pattern of types}:
only its flexible variables may be assigned. Leaf checks parallel the
unifier; in particular an ordinary pattern variable cannot match a
result-capable target variable. On ordinary product shapes match each pair
\emph{independently}, merging maps only when they agree on common domains.
A repeated source variable must therefore meet the same target twice.
For nested $\forall$, open the source flexibly, target rigidly and check
the same bijection/non-escape conditions. In predicate/instance comparisons
freshening apart matters: matching is not unification with an occurs check.
For instance $(a\arr a)\oneway(b\arr\Int)$ fails when $b$ is a distinct
unconstrained target. Symmetric unification would equate $b$ with $\Int$,
which is exactly what a dictionary lookup is not allowed to do.

\paragraph{The first-order fragment.}
Take finite well-kinded types with flexible variables, rigid leaves,
nominal constants, application, and closed keyword-free functions of fixed
positional arity and one ordinary output. Kinds are fixed; use normalized,
alias-free inputs and substitution ranges, with no packs, quantified types,
origins or result-capable variables. ``First-order'' here refers to the
term algebra, not to a restriction that every variable have kind $*$.

Use one mathematical tree for each type: read every unary function
$A\arr B$ as $(\mathsf{Arrow}\,A)\,B$. The constant $\mathsf{Arrow}$ and
its partial applications are allowed. A function of any other arity $n$
has a distinct tree constructor with $n$ inputs and one output.
For example, with $f:*\arr *$, both unifying $f\,a$ with
$\Int\arr\Int$ and matching $f\,a$ into it give
$[f\mapsto\mathsf{Arrow}\,\Int,\ a\mapsto\Int]$. Although neither input
mentions \code{Arrow} by name, the substitution needs its partial application.

\begin{theorem}[First-order solver]\label{thm:mgu}
For a well-sorted list of equations in this fragment, ordinary unification
from the empty substitution succeeds iff there is a finite kind-preserving
solution fixing the rigid constants. Its answer $S$ solves the equations,
and every solution $U$ factors as $V\circ S$ on the input variables.
For fresh-apart trees in the same fragment, one-way matching succeeds
iff the target is a substitution instance of the source; it does not
alter target variables. Equality and factorization are on the trees above.
\end{theorem}
\begin{proof}
Substitution commutes with this tree interpretation: the only change of
representation is that a newly saturated \code{Arrow} is stored as a unary
function, which denotes the same tree. When comparing an application with
a unary function, \src{mgu} and \src{match} expose the latter's arrow spine.
When both sides are functions they compare their slots directly, which
decomposes the same trees. No such spine exists for the other arities.

Maintain a solved substitution and a work list. Removing a reflexive equation
does not change the solution set. Distinct constants or incompatible shapes
have no common instance in a free term algebra. Equal application or function
heads are injective in every slot, so replacing their equation by the ordered
child equations is equivalent. For $a\eqs t$ with $a$ flexible, unequal
kinds cannot be repaired at this stage. A proper occurrence of $a$ in $t$
would make the finite tree $Ua$ a proper subtree of itself. Otherwise put
$B=[a\mapsto t]$. Every solution satisfies $Ua=Ut$, so $U=U\circ B$
on the variables; conversely a solution of the remaining $B$-transformed
equations composed with $B$ solves the old list. Orienting a flexible/rigid
pair in the other direction would invalidate this argument. Iterating the
factorization proves most-generality and the composition identity above proves
soundness. Termination follows lexicographically by the number of distinct
flexibles in the unsolved equations, then total unsolved tree size:
elimination removes a variable, decomposition or deletion reduces size
when the first count does not fall. These are exactly the recursive
decompositions and threaded calls of \src{mgu} and \src{unify-list}.
For matching, induct on the source tree: each variable occurrence prescribes
one target subtree; equal prescriptions on shared variables are necessary
and sufficient to take the union. A rigid leaf must agree literally.
Fresh-apartness makes the resulting map disjoint from target variables.
\end{proof}
\noindent The same elimination argument on the grammar $*,u,\kappa\arr\lambda$
applies to kind equations \emph{before} the final $*$ default. This theorem
does not say that term inference with defaults or branch result choices has
principal types.

\sources \src{substitutions.lisp}: \src{apply-substitution},
\src{compose-substitution-lists}, \src{merge-substitution-lists};
\src{unify.lisp}: \src{bind-variable}, \src{mgu}, \src{match}, \src{unary-function-type-as-arrow-type},
\src{mgu-forall-types}, \src{validate-forall-bijection-and-filter-subs}.

\section{Classes, evidence and improvement}\label{sec:classes}

\subsection{What proves a predicate?}
A class $C\,\bar a$ declares superclass fields $P$, method signatures,
and dependencies between positions of $\bar a$. A method with signature
$Q\Rightarrow t$ is available as a value with context
$[C\,\bar a]\mathbin{+\!\!+} Q$. Class parameters and method-local binders
are distinct. Superclass hierarchies are acyclic. A superclass field may
mention an extra variable absent from the subclass head: two occurrences of
the class do not thereby share the same hidden witness. At an explicit
binding or expected $\forall$, \src{instantiate-context-superclasses}
materializes hidden fields (also through intermediate superclasses).
These fields are evidence attached to a predicate occurrence, not extra
class arguments. Substitution carries the fields along; instantiating a
scheme for a new use does not copy the old dictionary's witnesses.
Improvement and superclass lookup must see the same fields. If
$\mathsf{Child}\,a$ supplies $\mathsf{Assoc}\,a\,h$ with $a\to h$,
adding a given $\mathsf{Assoc}\,a\,b$ can identify $h$ with $b$ in
either order. $\mathsf{Child}\,a$ alone does not name a caller-chosen $b$.

An instance is a named clause $i:\forall\bar b.\,P_i\Rightarrow h_i$, with
an optional \code{overlap} mark. Alpharenamings apart are implicit when
comparing heads. Say $h_1\succeq h_2$ (``at least as general'') when
$h_1\oneway h_2$ succeeds, matching class and arguments; $h_2$ is strictly
more specific if the converse fails. Mutual matching heads are replacements,
not a pair of distinct alternative candidates. A new, distinct head that
\emph{unifies} with an existing one requires \emph{both} instances to be
marked. For a wanted $p$:
\[
 M_\Sigma(p)=\{i\in\Sigma:h_i\succeq p\},\quad
 W_\Sigma(p)=\{i\in M_\Sigma(p):\text{no }j\in M_\Sigma(p)
                      \text{ is strictly more specific}\}.
\]
$\sel_\Sigma(p)=i$ iff $W_\Sigma(p)=\{i\}$ and either $\fv(p)=\varnothing$ or $i$ is
unmarked. Otherwise there is no stable choice; incomparable winners for
a ground $p$ signal ambiguity. Do \emph{not} try the most specific context
and backtrack to a weaker head when a prerequisite fails. A non-ground
constraint with a marked winner can be retained as a dictionary parameter.
Recorded instance choices and some optimizer assumptions about classes
without marked instances are validated when the environment changes.

Distinguish the \emph{declared} context $P_i$ from the compiled interface
$I_i=(h_i;Q_i)$, where $Q_i$ is an ordered vector of dictionary inputs.
Initially a new entry has the provisional vector $P_i$. Once all heads
of the unit are installed, \src{finalize-instance-signatures} computes
the vectors for the entire batch \emph{before} replacing any provisional
vector. A prerequisite with no stable instance is retained as a leaf
unless it leads back to the defining instance as described below
(an ambiguous ground selection is an error). For a selected
instance, expand its instantiated interface inputs in order: these come
from a provisional signature for this batch and from the saved signature
for an older entry. Duplicate leaves are removed by predicate equality
(the last occurrence survives). The result, including the head with its shared variables, is
saved as $I_i$; a one-element list $(h_i)$ means no dictionary inputs.
Absent signature metadata is an error requiring recompilation, not the
nullary case.

\paragraph{Residual prerequisites.}
Expansion records the selected instance, goal, depth and prerequisite
position at each frame, including the defining instance at depth zero.
The top-level prerequisites are numbered; a derived instance sets the
position to its child's index, while a non-derived instance inherits a
specified position (or uses its child's index if none is specified).
The root and superclass obligations have unspecified positions.
Goal keys keep variables of the original head and context fixed and compare
path-local variables up to renaming. Write $p\preceq_X q$ for the
implemented homeomorphic embedding with fixed variables $X$: the class
and arity agree and each earlier argument embeds in the corresponding
later argument. A type embeds by descending into a child or by matching
its root and recursively embedding corresponding children. Fixed variables
match only themselves at a variable leaf; a path-local variable can couple
only with another nonfixed variable. Functions respect their positional
and keyword shapes, and the $\forall$ case aligns compatible binders;
descending under a quantifier must not capture a variable of the earlier type.

On reselecting the same instance, first look for an equal goal key. Otherwise
look for an earlier embedded goal at a compatible prerequisite position
(equal positions, or either unspecified). Ordinary recursion unwinds to the
entry of that cycle or growth, but no farther than the first prerequisite
below the root. The predicate at that frame becomes a \emph{required input},
not an empty proof. Expansion also checks superclass paths for back-edges
but does not append their leaves to the instance's input vector;
the superclass check in \src{typecheck-instance} seeds the path with the
instance being checked, so it cannot justify its own missing superclass.

When no instance is stably selected at depth greater than one, expansion
also checks whether the defining head matches the goal and its instance
has a matching frame by the same key or embedding test. If so, it unwinds
to retain the prerequisite as above. This is cycle detection, \emph{not}
selection. In particular an open goal matching a marked head cannot use
this fallback to discharge an input, even if the path includes a derive.
In the \code{AuditTree} example in \src{tests/fundep-tests.lisp}, the
generic \code{CarrierEq} instance thus keeps its
$\mathsf{Eq}\,(\mathsf{AuditTree}\,a)$ input for comparing the contents;
a ground client can still select the more specific wrapper instance.

There is a distinct convention for a \emph{stably selected} structural back-edge.
\src{valid-derived-cycle-p} requires an equal key and a cycle containing a
compiler-generated instance. Walk the cycle from the newest frame back to
the matching frame. Start with the oldest derived frame's predicate (the
current goal if that frame is the matching frame), and reset the target
at each derived frame encountered. At each non-derived frame, every
argument of the current target must be a proper subterm of an argument
of that frame's predicate. An accepted structural back-edge
contributes no additional input; the defining instance need not itself be
derived. An embedded but unequal key does not close such a cycle. Thus
regular recursive derivations can close, whereas a nested datatype at a
larger parameter can retain a field constraint.

Write $\Sigma;G\vdash p\leadsto d$ when the available dictionaries
$G=[x_i:p_i]$ supply a proof term $d$. Its meaning is $d:\Dict(p)$. For
an ordered path $\pi$ of superclass fields, $\pi(x)$ projects them in
order (the empty path returns $x$). In \textsf{Given}, $E_G$ is the
identity except for local equations proved by supplied pattern dictionaries
(Section~\ref{sec:patterns}); it is not obtained by matching the request.
In \textsf{Instance}, $(h;q_1,\ldots,q_n)$ is one fresh copy of the
selected instance's stored interface. The rules behind \src{entail} and
\src{resolve-dict} are
\[
 \rrule{Given}{(x:q)\in G\quad q\xrightarrow{\pi}_{\Sigma}r\quad E_G r=E_G p}
                {\Sigma;G\vdash_H p\leadsto\pi(x)}
\]
\[
 \rrule{Instance}{\sel_\Sigma(p)=i\quad h\oneway p\leadsto M\quad
       \mathrm{open}_{H}(i,p)\quad
       \Sigma;G\vdash_{H,(i,p)}\theta q_j\leadsto d_j\ (1\le j\le n)}
       {\Sigma;G\vdash_H p\leadsto i(\bar d)} .
\]
The head match gives $M$. A variable shared by two inputs but absent
from $h$ is fresh \emph{once per call}, not once per input.
\src{by-inst} takes $\theta=M$; \src{resolve-static-dict} also solves
FDs on $M\bar q$, obtaining $F$, and uses $\theta=F\circ M$.
The typing argument requires $\theta h=p$; an undetermined prerequisite
variable is not a free choice of dictionary. When $n=0$, $i(\bar d)$
denotes the dictionary value $i$, not a nullary call. Translation tries
supplied dictionaries and their fields before static instances or
registered synthesis. For example, when no ordinary instance is selected
it can synthesize a prerequisite-free $\mathsf{RuntimeRepr}$ dictionary
for a function type.

The annotation $H$ is the active instance/goal path (empty when omitted).
\textsf{Given} is available at every depth and is tried first.
$\mathrm{open}_H$ always blocks a previously selected instance at the
same goal key. A distinct goal $p$ at the same instance is a candidate
for the growth cutoff if some earlier $p_0\preceq_X p$ and the height
of $p$ exceeds \emph{both} that of $p_0$ and the search allowance.
This second cutoff applies only if the cycle shows no progress as
defined below. Height counts the longest type-tree path (a leaf has
height one); a predicate takes the maximum of its argument heights.

For an ordinary instance of class $C$, let $R_\Sigma(C)$ be the least
set containing $C$ and closed under direct superclasses and the classes
mentioned in the declared prerequisites \emph{or} stored dictionary inputs
of each instance of a class in the set. Let $m_\Sigma(C)$ be the maximum
head height over these classes, or zero if there are no heads. The allowance
is $h_0+m_\Sigma(C)$, not a bound from $C$'s own heads alone; a synthesized
step uses zero for the head contribution. For \src{entail}, $h_0$ is the
maximum height of the root goal, givens and their superclass closure;
$X$ comes from the root goal, given predicates and any initial stack.
Static instance resolution uses the root goal and the superclass closure
of its dictionary context for both $h_0$ and $X$. Definition-time
retention uses neither this height allowance nor the following test.

\paragraph{Progress toward a repeated-variable head.}
For a growth-cutoff candidate at instance $i$, inspect the current step
and the intervening steps since the \emph{most recent} prior selection
of $i$, excluding that prior step. Any one of these frames can permit
search to continue. To test an ordinary frame $(j,q)$:
\begin{enumerate}
\item For every \emph{other} head of $q$'s class, compare the argument
      subtrees of $q$ at the first and each later occurrence of a repeated
      head variable. Follow child positions; no head match is required
      to take this measurement. Record $\infty$ when an operand is absent.
      The disagreement is zero for identical nodes or equal leaves. For
      nonleaves of the same internal node class and child count, it is the
      maximum disagreement of corresponding children; otherwise it is the
      larger operand height. If none of these \emph{head} distances is
      finite, the frame cannot report progress. A finite distance improves
      on $\infty$; even distance zero is not a proof of type equality.
\item The \emph{carried} variables are those occurring both in $j$'s
      head and in its stored dictionary inputs. Match that head into $q$.
      For each pair of carried variables with instantiated heights $u,v$,
      record both $\max(0,u-v)$ and $\max(0,v-u)$. Keeping directions
      separate distinguishes passing the other input from moving farther
      away on the same side.
\end{enumerate}
For each earlier visit to $j$ in its full preceding history, first compare
the head distances. Only if none decreased for that visit must a carried
gap decrease instead. Equivalently, once a head distance is observable,
\emph{each} earlier visit needs \emph{some} strictly smaller coordinate
in the combined vector $V_j(q)$; the coordinate may differ between
visits. With no earlier visit the comparison is vacuous. Synthesized
frames never report progress. No measurement selects a head or supplies
a dictionary.

\begin{lemma}[Finite progress reports on a fixed path]
Fix the signature and an append-only path with unchanged observations
for earlier visits (in particular, a path of ground goals). An instance
can report progress at only finitely many of its own visits.
\end{lemma}
\begin{proof}
Its vector has fixed finite dimension and entries in
$\mathbb N\cup\{\infty\}$. An infinite sequence in one coordinate
has an infinite nondecreasing subsequence: either a value recurs
infinitely often or only finitely many entries lie below any given
natural bound and $\infty$ occurs only finitely often. In the latter
case choose successively larger natural entries. Thin an infinite
sequence of reporting visits this way once per coordinate. Two visits
in the resulting subsequence have
no coordinate strictly smaller at the later visit, contrary to the test.
\end{proof}
The lemma is about fixed observations on one path; it is not a
termination theorem for the compiler or for histories being resubstituted.

Methods of an instance are checked at the class signatures specialized to
its head; the instance context supplies dictionaries but does not permit
the implementation to narrow the public head. Required superclass
dictionaries must be available without further specializing the variables
fixed by the instance and its givens. The resulting dictionary holds the
superclass fields and method implementations. \src{typecheck-instance}
uses $Q_i$ as the method context, and \src{translate-instance} emits the
corresponding parameters. A later client can choose \emph{how to supply}
these dictionaries in its own environment, not change their types or order.
Derivation generates instance declarations; exceptions and resumptions get
sealed marker instances, and runtime representations have special rules.

\begin{theorem}[Selection and evidence, in a fixed signature]\label{thm:evidence}
Let $p$ and the instance heads of its class belong to the first-order
fragment of Theorem~\ref{thm:mgu}, and suppose those heads satisfy the
overlap admission rule. An unmarked selected head for $p$ remains the
unique selected head at every ground substitution instance $Up$ in that
\emph{same} $\Sigma$. Furthermore, for finite successful evidence search
with well-typed given fields, dictionary constructors typed at their
stored $Q_i\Rightarrow h_i$, and correctly typed synthesized dictionaries,
the resulting evidence has type $\Dict(p)$ up to valid local conversion.
Here a conversion in \textsf{Given} is valid if the projected and requested
predicates agree at every ground valuation satisfying the local givens
and index equations. Section~\ref{sec:fd} establishes this premise
for first-order fields satisfying the declared FDs.
Any extra FD instantiation must respect the scope of
prerequisite-only variables and preserve the goal and given types.
\end{theorem}
\begin{proof}
If $Mh_i=p$, composition gives $(U\circ M)h_i=Up$. Any other head matching
$Up$ would have a common instance with $h_i$; the fresh-apart assignments
combine to a unifier. Admission forbids this since $i$ is unmarked, unless
the heads are replacements, in which case only one is present. For the
second assertion induct on search height. A given is either the correctly
typed input dictionary or a chain of typed record projections; induction
on path length gives the projected field type $\Dict(r)$. If equality is
not literal, the local conversion premise licenses $\Dict(p)$ at the
same ground instantiations. At an instance step,
$\theta h=p$ specializes the stored constructor's result to $\Dict(p)$.
By induction argument $j$ has type $\Dict(\theta q_j)$; application in
the stored order therefore has the required type (or take the dictionary
value directly when $n=0$). Synthesis is the stated base case.
The growth guard contributes no proof term. Whether it permits the step
by height or by progress, the same selected constructor still requires
all of its stored inputs.
\end{proof}
This proves neither equality of the \emph{values} of dictionaries from
different paths nor termination. In particular, cutting cycles in
\src{expand-context} (used for recursive derived instances) is \emph{not}
induction on a finite proof in the Horn theory of the original clauses:
the raw clause $p\Rightarrow p$ alone has no such proof of $p$. To read
the evidence argument as ordinary logical entailment from the raw clauses,
require finite acyclic context expansion and correctly instantiated
witnesses, including for older interfaces. To include generated structural
back-edges in a semantic soundness argument one must justify the recursive
dictionary construction separately; the finite Horn induction does not do
that. The convergence examples below show finite paths that the checker
does find; these arguments do not establish completeness of proof search.

\paragraph{Linking corollary.}
Let a provider emit $i$ at $I_i$ and a client select it in a later
environment. For an unchanged instance, or a replacement with the same
compiled interface up to a single binder renaming, the client specializes
the stored inputs under the \emph{same} instantiation as the head.
Indeed, loading the provider restores $I_i$. Compose the freshening of
that signature (and any binder renaming) with $\theta$ for this call.
The resulting map specializes the stored head to $p$ and the $j$th
parameter to the predicate proved for the $j$th argument, in order.
The evidence induction therefore applies even when a new instance is
used in one of those argument proofs. This corollary has the typing
hypotheses of Theorem~\ref{thm:evidence}; it assumes compatible class
layouts and does not assert that method \emph{values} never change.

For recorded selections, validation reselects an instance for the recorded
predicate and compares the overlap and derivation flags, constructor symbol and joint
ordered signature, including classes, arities, kinds and type-variable
sharing. Unit compilation defers checking changed classes until the batch
is finalized; FASL replay checks the stored interfaces before installing
the resulting environment. Checking the head
alone, the raw source context, or the parameter count alone would not
suffice. Optimizer reconstruction likewise uses the stored signature.

\begin{example}[A fixed input supplied by new evidence]\label{ex:dict}
Let $W$ be a unary datatype, and $N,O$ unary classes. In $\Sigma_0$ install
$N\,\Int$ and $i:\forall a.\,N\,(W\,a)\Rightarrow O\,a$. No instance
matches $N\,(W\,a)$. Compile $i$ now. In $\Sigma_1$ add
$j:\forall b.\,N\,b\Rightarrow N\,(W\,b)$. It does not overlap
$N\,\Int$. At $O\,\Int$ the two environments select the same $i$
with the same stored parameter $N\,(W\,\Int)$. What changes is the
availability of an argument, not the interface:
\[
\begin{aligned}
 d &: \Dict(N\,\Int),\qquad j(d):\Dict(N\,(W\,\Int)),\\
 i\bigl(j(d)\bigr)&:\Dict(O\,\Int).
\end{aligned}
\]
Re-expanding $i$ at the use site would instead pass $d$ to $i$: the arity
would agree but the argument type would not. Recompiling $i$ in
$\Sigma_1$ is a different operation: its freshly computed interface can
change and previously recorded callers must then be revalidated.
\end{example}

\subsection{What determines a type?}\label{sec:fd}
For $C$ with a functional dependency $X\to Y$ (sets of parameter positions),
the intended assertion about a ground dictionary relation $R_C$ is
\[
 \bar u,\bar v\in R_C,\quad \bar u|_X=\bar v|_X
       \quad\Longrightarrow\quad \bar u|_Y=\bar v|_Y.       \tag{FD}
\]
It is not a dictionary constructor. Improvement changes types; the
resulting predicate still needs its own evidence. For instance
$C\,a\,b,\ C\,a\,c$ with $1\to2$ gives $b\eqs c$, not $a\eqs b$,
and $C\,a\,b,\ C\,d\,c$ alone does not give $a\eqs d$.

For each predicate in a context and each inherited superclass predicate,
collect pairs
$(\fv(\bar t|_X),\fv(\bar t|_Y))$. For a finite such set $F$, let
\[
 \clo_F(A)=\mu Z.\,\bigl(A\cup Z\cup
                  \bigcup\{V:(U,V)\in F,\ U\subseteq Z\}\bigr)
\]
(start at $A$). In a declared qualified type all predicate variables must
lie in $\clo_F(\fv(t))$. A method must identify all class head parameters
from the head parameters occurring in its type using class and superclass
dependencies. A class's own dependencies may name only its head parameters.
A declaration is also checked for actual restriction by improvement; a
context reducible merely by entailment can receive a warning.

\paragraph{Instance admission.}
For every dependency of the head require coverage:
$\fv(h_i|_Y)\subseteq\clo_{F(P_i)}(\fv(h_i|_X))$. Write $(l,r)$ for
the head's determinant and dependent projections. Compare the new clause
with \emph{every other definition} of the same class, freshening each
head together with its declared prerequisites. Skip the equivalent head
being replaced. The improvement table can coalesce projected rows with
different prerequisites; it is not the list of definitions to compare.
For two such clauses $(l,r,P)$ and $(l',r',P')$:
\begin{enumerate}
\item A failed determinant unification ends the comparison (in the
      first-order fragment, the two clauses then have no shared input).
      Otherwise let $S=\mgu(l,l')$, $u=Sl$, $v=Sr$, and $v'=Sr'$.
      If $v=v'$ pointwise, agreement is immediate.
\item Otherwise the two projected rows $(u,v)$ and $(u,v')$ must
      one-way match in \emph{both} directions. The separate improvement
      index can use rows without their contexts; it may not promise a
      more specific output shape for one of these rows than the other.
\item Take only $S(P\mathbin{+\!\!+}P')$ and their superclasses as givens.
      Repeatedly compare same-class predicates with already equal
      determinants and unify their dependents, until the substitution stops
      changing. A conflict, or failure to obtain $Tv=Tv'$ at that fixed
      point $T$, rejects the new clause. Neither head is a given in this
      proof: unifying $v$ with $v'$ simply because the heads declare an
      FD would assume the result. Even contradictory prerequisites do not
      count as proof by vacuity. No lookup of today's ground instances is
      a premise.
\end{enumerate}
The checker also tests a \emph{single} head and its context for internal
conflicts. It maintains the projected improvement index separately.
Admission is conservative: currently unsatisfied contexts do not certify
disjointness under future extensions.

\paragraph{Solving.}
After substituting the known equations, expand superclasses (reusing
given expansions when supplied). \src{solve-fundeps} iterates:
\begin{enumerate}
\item Compare predicates of the same class: for each FD with already
      \emph{equal} determinant vectors, unify its dependent vectors,
      signalling a conflict on failure.
\item A stable instance match can expose its compiled inputs
      $Q_i$, freshly instantiated for further improvement. Freshen the
      head \emph{with} the context so hidden context variables are shared
      within that occurrence only. A generated predicate remembers its
      instance path. On a repeated selection the solver substitutes into
      that history and uses the same repeated-key, growth and progress
      tests as evidence search.
      Here $h_0$ and $X$ are the maximum height and free variables of the
      \emph{original input predicates} under the round's substitution,
      not of all generated prerequisites. A blocked instance is not unfolded.
\item Without an unblocked match, first try instance-head improvement.
      Even without an FD, a unique unmarked \emph{unifying} instance can
      constrain repeated variables \emph{inside one structural argument}.
      Freshen its argument patterns separately; its linearization
      (new variable per occurrence) must already match the wanted shape.
      Keep only wanted-variable refinements caused by repetition, with no
      surviving fresh instance variables. E.g.\ a sole $C\,(D\,x\,x)$
      can relate $a,b$ in $C\,(D\,a\,b)$; a sole $C\,\Int$ does not
      default $C\,a$ to it. This is signature-relative, not (FD).
      A successful improvement ends this scan so that the next round
      compares determinants with the new substitution.
\item Otherwise, for a class with FDs, try its determinant table:
      match a fresh entry's determinant into the wanted determinant,
      unify the resulting dependents, and keep an improvement only if
      it makes the wanted predicate strictly more specific. Nonapplicable
      entries, including failed unifications here, are skipped.
\end{enumerate}
The iteration is bounded at 1000 rounds; the solver distinguishes continued
specialization from a failed fixed point on exhausting that budget. When
comparing declared/given to wanted predicates, \src{fundep-entail} can
further align differing \emph{determined} arguments. After superclass
expansion it takes as known positions those mentioning a known variable,
closes those positions under the class FDs, finds a same-class given with
the same known arguments, and one-way matches determined wanted arguments
to the corresponding given arguments. The checker iterates this with a
growing known-variable set and checks entailment afterwards.

\begin{lemma}[When pairwise improvement is sound]\label{lem:fd}
Use the trees of Theorem~\ref{thm:mgu} and assume relations satisfying
(FD). For an ordinary unification step, valuations fix its rigid leaves.
For the branch-local computation of $E_G$, the skolems opened by
\src{unify-given} are parameters of the ground valuation, not distinct
ground constants; the valuation must satisfy the supplied predicates,
their actual superclass fields and the branch equations. If $U$
is a ground solution of two same-class predicates and their determinants
already agree under $S$ through which $U$ factors, unifying their dependents
under $S$ loses no such solution. For the collected dependencies, if two
satisfying assignments agree on $A$, they agree on $\clo_F(A)$, assuming
they satisfy the same expanded predicate occurrences (including their
superclass witnesses).
\end{lemma}
\begin{proof}
The determinants remain equal under $U$. By (FD) the instantiated
dependents are equal, so Theorem~\ref{thm:mgu} factors $U$ through the
ordinary unifier. For the local step, assign each fresh flexible copy
the ground type of the skolem it opens. The same factorization now
applies. Closing a surviving representative back to an original name
preserves that valuation: names unified with the same representative
already have the same ground type. This is not a binding in the caller's
inference substitution. For closure, agreement on all variables of a
determinant vector implies agreement on the instantiated vector. By (FD)
the instantiated dependent trees agree. Every occurrence of a variable
in those trees is at a fixed path in a free term tree; equality at that
path gives agreement on the variable. Induct on closure rounds. Neither
claim establishes (FD) for a particular set of declarations.
\end{proof}
In particular, under these first-order hypotheses each step producing
$E_G$ preserves every admissible ground valuation. Hence $E_G r=E_G p$
implies $Ur=Up$ at each such valuation. This supplies the local conversion
premise of Theorem~\ref{thm:evidence} without assuming that two dictionary
values with the same type behave identically.

\paragraph{What admission proves for two clauses.}
Suppose two distinct admitted heads have ground uses with equal inputs
and with both declared contexts true. In the first-order fragment their
joint valuation factors through $S$. If $v=v'$ the outputs already agree.
Otherwise assume the prerequisite relations used in step 3 satisfy (FD)
and their superclass fields are valid. Extend the valuation with witnesses
for the fresh context-only and hidden superclass variables. By
Lemma~\ref{lem:fd}, every improvement step preserves a solution extending
the original head valuation. Hence equality
$Tv=Tv'$ at acceptance implies equality of the two ground outputs.
The mutual row-shape check further restricts admission; it is not the
proof of that equality. For two uses of \emph{one} clause, agreement on
the input trees gives agreement on their variables; coverage and the
closure part of the lemma give agreement on the output variables under
the same prerequisite hypotheses. If a prerequisite relies on the head
class's own FD, that hypothesis cannot establish functionality from
scratch. The argument is noncircular along well-founded prerequisite
proofs. In particular, retaining a recursive prerequisite as an input
does not establish it without a supplied dictionary.

\begin{example}[A rejected pair and a shared-prerequisite control]\label{ex:fd}
In a fresh package using Coalton, the following program is accepted:
\begin{verbatim}
(define-class (LKey :k :v (:k -> :v)))
(define-class (RKey :k :v (:k -> :v)))
(define-class (TaggedKey :k :v :tag (:k -> :v)))
(define-instance (LKey Unit Integer))
(define-instance (RKey Unit String))
(define-instance (LKey Unit :v => TaggedKey Unit :v Boolean))
(define-instance (LKey Unit :w => TaggedKey Unit :w Char))
\end{verbatim}
The full heads are disjoint ($\Bool\ne\mathsf{Char}$), but their determinant
is $\Unit$ in both cases. Coverage alone permits the two output variables.
Agreement follows because the two \code{LKey} prerequisites have the same
determinant, so their declared FD equates the outputs. Replace only the
last prerequisite's \code{LKey} with \code{RKey}: neither given now equates
those outputs, so admission rejects whichever definition comes second.
\end{example}

\begin{example}[Retention and reachable base instances]\label{ex:fd-cycle}
Let $C\,k\,v$ be a method-free class with $k\to v$. The single clause
$C\,\Unit\,v\Rightarrow C\,\Unit\,v$ passes coverage through its own
prerequisite and has no other head to check for disagreement. Its compiled
input vector is $[C\,\Unit\,v]$. In given-free search for
$C\,\Unit\,\Int$, selecting the clause requires the identical goal at the
same instance, which is blocked. The independent $\Str$ query likewise fails.
A supplied matching dictionary is still evidence. The same retention
applies to a two-class hand-written cycle.

For $G\,(\List\,a)\Rightarrow G\,a$, the next required goal grows. Its
stored input is $G\,(\List\,a)$, not the result of an infinite unfolding.
With no givens and no other instances the use-site guard eventually blocks;
giving $G\,(\List^3\,\Int)$ instead supplies a finite path for $G\,\Int$.
Conversely $D\,a\Rightarrow D\,(\List\,a)$ with base $D\,\Int$ proves
$D\,(\List^5\,\Int)$ by descent, but not $D\,(\List^5\,\Str)$.

Let $\mathsf{WalkA}\,a$ require $\mathsf{WalkB}\,(\List\,a)$ and
$\mathsf{WalkB}\,a$ require $\mathsf{WalkA}\,a$, with both general instances
marked for overlap. Add a marked, more-specific, input-free instance for
$\mathsf{WalkB}\,(\List^3\,\Int)$.
Writing $L=\List$, the finalized inputs give the ground path
\[
\begin{aligned}
 \mathsf{WalkA}\,\Int &\leadsto \mathsf{WalkB}\,(L\Int)
                        \leadsto \mathsf{WalkA}\,(L\Int)\\
 &\leadsto \mathsf{WalkB}\,(L^2\Int) \leadsto \mathsf{WalkA}\,(L^2\Int)
                        \leadsto \mathsf{WalkB}\,(L^3\Int).
\end{aligned}
\]
Here the arrows mean ``requires''. More generally, if the ground base is
$k\ge1$ lists deep, the root height is one and both classes can reach the
base head of height $k+1$. The allowance is $k+2$; no goal on the path
exceeds $k+1$. Thus the height guard does not prevent this proof. Removing
the base or changing its leaf to an unrelated type leaves the goal unproved.
\end{example}

\begin{example}[Growing arguments that converge]\label{ex:fd-height}
Give $\mathsf{MeetDepth}\,a\,b$ a marked instance requiring
$\mathsf{MeetDepth}\,(L^2a)\,(L b)$, where $L=\List$, and a more specific
marked instance for $\mathsf{MeetDepth}\,s\,s$ with no inputs. Starting
at $\mathsf{MeetDepth}\,\Int\,(L^3\Int)$, write each pair of argument
types as their list depths over $\Int$. The stored inputs give
\[
 (0,3)\;\leadsto\;(2,4)\;\leadsto\;(4,5)\;\leadsto\;(6,6).
\]
Both heads have height one and the root goal has height four. At $(4,5)$
the goal height is six, above the allowance of five, and an earlier goal
embeds in it at the same general instance. But the disagreement of the two operands
at the first three goals is respectively $4,3,2$. Search continues; at
$(6,6)$ ordinary selection picks the more specific, input-free instance.
When the base method returns zero and the general method adds one, the
ground clients at initial gaps three and four return three and four.
Changing only the leaf of the second operand to $\Str$ gives no such base.

Progress need not come from the compared operands alone. Let a marked
$\mathsf{Stride}\,a\,b\,c$ require $\mathsf{Stride}\,(Lc)\,(L^2b)\,(L^2a)$,
with an input-free marked base at $\mathsf{Stride}\,s\,s\,c$. The list depths follow
\[
 (2,0,2)\;\leadsto\;(3,2,4)\;\leadsto\;(5,4,5)\;\leadsto\;(6,6,7).
\]
At $(5,4,5)$ the first two arguments are still one layer apart, just as
at $(3,2,4)$. But the height of $c$ exceeds that of $b$ by one rather
than two, and $c$ feeds the first argument of the next prerequisite.
Against the initial visit the repeated-variable disagreement itself has
decreased. The step is therefore allowed and the next goal selects the
base. By contrast, growing two unequal arguments in lockstep does not
make them equal. Continuing a search and successfully discharging its
last prerequisite are distinct events.
\end{example}

\paragraph{Why the revision matters.}
On \code{deb300be}, the client in Example~\ref{ex:dict} passed $d$ rather
than $j(d)$, and the split-prerequisite variant in Example~\ref{ex:fd}
entailed both $\mathsf{TaggedKey}\ \Unit\ \Int\ \Bool$ and
$\mathsf{TaggedKey}\ \Unit\ \Str\ \mathsf{Char}$ separately. On
\code{cd7cb56c}, the hand-written cycles in Example~\ref{ex:fd-cycle}
had zero-input signatures and entailed both ground results. The
class-local bound on \code{09c5e3fe} missed the \code{WalkA} proofs at
depths three and four; the reachable-head bound on \code{21c6d7c8}
still stopped the \code{MeetDepth} example at $(4,5)$. These are
contrasts with the current rules, not failures on the pinned revision.

\sources \src{define-class.lisp}: \src{infer-class-kinds};
\src{define-instance.lisp}: \src{define-instance-in-environment},
\src{typecheck-instance}; \src{environment.lisp}: \src{add-instance},
\src{lookup-class-instance-in-list}, \src{expand-context}, \src{finalize-instance-signatures},
\src{instantiate-instance-context}, \src{instance-signature=},
\src{predicate-embeds-p}, \src{valid-derived-cycle-p}, \src{max-class-instance-head-height},
\src{type-disagreement-height}, \src{head-equality-distances}, \src{instance-input-progress-p},
\src{instance-equality-progress-p}, \src{instance-recursion-blocked-p},
\src{instantiate-context-superclasses}, \src{update-instance-fundeps},
\src{check-fundep-agreement}, \src{context-fundep-substitutions}, \src{improve-predicate-fundeps},
\src{solve-fundeps}; \src{context-reduction.lisp}: \src{entail},
\src{by-inst}, \src{fundep-entail};
\src{src/codegen/resolve-instance.lisp}: \src{resolve-dict}, \src{resolve-static-dict};
\src{src/codegen/translate-instance.lisp}: \src{translate-instance}.

\section{Checking expressions}\label{sec:terms}

Write $\Gamma(x)=\sigma$ for lexical then global scheme lookup. $\Delta$
also tracks source type names currently in scope. $G$ contains available
predicates, $E$ the local type equations supplied by patterns, and $B$ the
lexical return targets. Abbreviate this data with $\Sigma$ by $\mathcal C$.
The main checking judgment is
\[
 \mathcal C;S\vdash e\chk t
      \ \Longrightarrow\ (u,W,A,e',S').
\]
It checks a parser expression $e$ against expected type $t$, returns its
inferred type $u$, wanted predicates $W$, delayed accessor obligations $A$,
and typed syntax $e'$ with updated substitution $S'$. At ordinary value
nodes $S'u$ agrees with $S't$, using the visible interface when there is a
coercion. An effect-only \emph{binding statement} need not even have a
meaningful $u$. Synthesis is a call with a fresh expected $a^v$, or $a^r$
where any result sequence is permitted. It is \emph{not} an unrestricted
search for a polytype. The short rules below omit unchanged data, carry the
substitutions left-to-right through the premises, concatenate obligations,
and attach the solved type to the corresponding syntax. In the expression
rules, $e\chk t$ is a check against an expectation; $e:t$ is a typing
premise or conclusion. Only $\vdash_0$ denotes the declarative core
judgment of Section~\ref{sec:safety}. $U_S(u,t)$ denotes
the ordinary equation step of Section~\ref{sec:unification}. Sections below
specify the cases that open a local scope instead of threading everything
outward.

\subsection{Values, functions and applications}
\[
 \rrule{Var}{\operatorname{inst}(\Gamma(x))=(P,u)
                  \quad S'=U_S(u,t)}
            {\mathcal C;S\vdash x\chk t
              \Longrightarrow(S'u,S'P,[],x^{S'P:S'u},S')}
\]
This is the ordinary lookup rule; constructors, methods and recursive
references use the scheme under their name. The superscript records the
predicate list and occurrence type used by translation. Integer literals
$n$ have the same rule with $(P,u)=([\mathsf{Num}\,a],a)$, $a$ fresh:
conversion is through a dictionary, not an initial guess of $\Int$.
Characters, strings, ratios and the two floating-point literal formats
have fixed types $\mathsf{Char},\Str,\mathsf{Fraction},\mathsf{F32},
\mathsf{F64}$ respectively, with no wanteds. Integer \emph{patterns}
also demand $\mathsf{Num}$; a literal test is not constructor-tag evidence.

\paragraph{Abstraction.}
For $(\mathsf{fn}\ \bar x.\,e)$ with ordinary variable parameters, take
the positional types $\bar u$ and result $r$ from an expected function
type when present, checking arity. Otherwise choose fresh single-value
inputs and a fresh result-capable $r$. With empty keyword set the rule is
\[
 \rrule{Fn}{\Gamma[\bar x:\bar u]\vdash e\chk r
               \quad \Fn(\bar u;\varnothing;\mathsf{closed};\pack^{-1}(r))
                      \eqs t}
            {\Gamma\vdash(\mathsf{fn}\ \bar x.\,e)\chk t}.
\]
Here and below a short $\Gamma$ judgment inherits the rest of $\mathcal C$.
There is no generalization of an inferred parameter. A parameter with a
known $\forall$ type binds the corresponding scheme instead. Patterns in
parameters are checked together (Section~\ref{sec:patterns}). With keyword
parameters, check each default at its expected keyword type or a fresh
ordinary variable, \emph{before} introducing its binder; later defaults see
earlier binders. The elaborated prefix chooses the supplied value or
evaluates the default, based on a presence flag. An abstraction's own
interface is closed.

\paragraph{Application.}
The meaning rule for a known closed function is
\[
 \rrule{Call}{f:\Fn(\bar u;K;\mathsf{closed};\bar r)
        \quad \bar e:\bar u \quad J\subseteq\dom K
        \quad(e_k:K(k))_{k\in J}}
       {f(\bar e;\,{:}k=e_k\ (k\in J)):\pack(\bar r)} .
\]
The vector premise requires \emph{exact} positional arity: this rule
does not curry, partially apply, or discard extra outputs. The checker
implements it as follows:
\begin{enumerate}
\item Infer the operator at a fresh ordinary variable. Its fresh
instantiation's result may be unified speculatively with the expected result
as a hint to its arguments, only when the hint binds exclusively variables
fresh since starting the operator. The final result check is still required.
\item If the operator is still a variable, unify it with a function of the
actual positional arity, fresh inputs and supplied keyword set, and outputs
from the expected result (or the direct \code{values} binding context).
The guessed function is open with a fresh origin iff keywords were supplied.
\item Check non-$\forall$ positional arguments first, then the remaining
polymorphic positions, then the supplied keyword arguments, putting the
typed terms back at their original positions. This is \emph{checking}
order, not a change in argument evaluation order. A missing label is an
error in a closed interface. For an open interface, extend it with the
label at a fresh type and a new origin, redirect the old origin to the
extension, and update an eligible unresolved local binding.
\item If the result is a bare $\forall$ in an ordinary context, instantiate
it, add the new wanteds, and account for their dictionary applications.
Unify the effective result with $t$. A temporary binds an expression whose
result must be consumed with dictionaries; instantiation does not duplicate
its evaluation.
\end{enumerate}
A variable or abstraction used as a function \emph{value} at an expected
function interface may hide optional keywords. For each expected label,
the actual interface must contain it, or be open and extendible. Project
both to those labels, close the projections, then unify their slots and
outputs. Thus a closed value accepting both \code{:x} and \code{:y} may be
used via a view exposing only \code{:x}, not conversely. This specific
coercion does not implement general function subtyping or row polymorphism.

\paragraph{Fields and packs.}
The accessor \code{.field} produces a unary $a\arr b$ with the obligation
$a.\mathit{field}=b$. If the head of $Sa$ is unknown, defer it. Otherwise
it must identify a struct with this field: freshen its parameter vector,
one-way match its head into $Sa$, and unify the specialized field type with
$Sb$. A bare unconstrained polymorphic field is instantiated here; a
qualified polymorphic field requires pattern matching instead. Iterate
accessor solving while at least one is solved; unresolved accessors at the
binding check are ambiguous. There is no row constraint that can be
generalized in their place.

For \code{(values} $\bar e$\code{)}, check each $e_j$ against a separate
fresh \emph{ordinary} $a_j$; unify $\pack(\bar a)$ with $t$.
A \code{let (values} $\bar x$\code{) =} binding requires exactly the
number of components named, each binder a variable or wildcard. A named
ordinary local value requires one result; an explicit discard
\code{let \_ =} can discard any result sequence. In a sequencing body the
nonfinal expressions have fresh result-capable expectations, the last
expression inherits $t$. Analysis may diagnose discarded non-void results
even when this equation check succeeds.

\begin{example}[The output slot in a higher-order call]
The quantifier in this declaration belongs to the caller, not the callback:
\begin{verbatim}
(declare call-results
  (forall ((:r Values)) (Void -> :r) -> :r))
(define (call-results f) (f))
\end{verbatim}
At a use with a nullary argument returning two results $\Int,\Str$, the
fresh $a^r$ for \code{:r} unifies with $\Pack([\Int,\Str])$ when checking
$f$'s output. The outer singleton output $[a^r]$ normalizes to
$[\Int,\Str]$. Neither binding the call to one ordinary variable nor
returning two values where a single $\mathsf{Tuple}\ \Int\ \Str$ is expected
uses an implicit conversion.
\end{example}

\sources \src{define.lisp}: \src{infer-expression-type}
(\src{node-variable}, \src{node-abstraction}, \src{node-application},
\src{node-values}, \src{node-values-bind}), \src{infer-application-rator},
\src{coerce-function-value-type}, \src{infer-keyword-params};
\src{accessor.lisp}: \src{solve-accessor}.

\subsection{Bindings and generalization}\label{sec:gen}
Local binding groups and top-level definitions use the following procedure.
First put explicit declarations in scope as schemes; infer the implicit
strongly connected components (SCCs) in dependency order; finally check
the declared definitions against their schemes. Dependencies include
keyword defaults. Sequential \code{let*} infers bindings successively;
the statement \code{let} with a pattern binds the checked pattern, not a
freshly generalized scheme at every variable in the pattern.

Within an implicit SCC, install mutually visible monomorphic placeholders
(an outer \code{the} on an initializer can supply its actual annotation
instead); infer all definitions, sharing equations. Collect their types
$\bar u$, wanteds $W$ and accessors $A$. Solve FDs before accessors,
then solve/reduce/default the context. In detail:
\begin{itemize}
\item $\operatorname{red}(W)$ applies current substitutions, removes
predicates entailed with empty givens, and removes duplicate or
superclass-implied predicates using the others. Builder defaulting can
choose $\mathsf{Seq}\,a$ or $\mathsf{Seq}\,(\mathsf{Tuple}\,k\,v)$ for
an unprotected unknown collection head, by unifying the \emph{whole}
collection type, solving FDs and exposing the builder's prerequisites.
\item Start the fixed variables $F$ with variables free in substituted
pre-group environment bindings, scoped type variables and relevant open
origins. Unresolved branch result choices and retained proof obligations
also protect variables from local quantification (Section~\ref{sec:patterns}).
Let $L$ be variables/origins of the inferred types outside $F$. Under the
relevant substitutions enlarge $L$ by FD closure minus $F$, then $F$ by
FD closure minus the enlarged $L$. In the reduced context split off
$D=\{p:\fv(p)\subseteq F\}$ (defer to caller); the rest $P$ is retained
provisionally.
\item A predicate variable not known from $F\cup L$ must be defaultable or
inference reports ambiguity. Numeric candidates for $a$ require that all
predicates mentioning $a$ have no \emph{other} variables, at least one has
a recognized numeric class, and all hold after the proposed substitution.
Choose the first valid candidate in $[\Int,\mathsf{F64},\mathsf{F32}]$.
Remove these defaulted predicates from $P$. At top level, also try defaulting
all $D$ (and all $P$ in a restricted group) and reduce again.
\end{itemize}

\paragraph{Relaxed value restriction.}
The syntactic predicate $N(e)$ holds for variables, literals, lambdas and
\code{type-of}; \code{the} preserves it, as does a single-expression
\code{unsafe} wrapper. A constructor applied to its \emph{exact} number
of positional arguments with all $N(e_i)$ is nonexpansive. An ordinary
function call is not, even a call to a function known to a human to be pure.
A function binding counts as nonexpansive; other bindings need a
nonexpansive initializer.

Variance signs are $+,-,\pm$ with $\bot$ for ``not observed''. Join
collects occurrences ($+\sqcup-=\pm$); an input, including a keyword
input, flips sign, an output/pack component preserves it, a nominal
argument composes with its stored variance, and an invariant or
unknown-head argument is invariant. Free variables in a nested $\forall$
are conservatively invariant. Compute variances for a mutually recursive
type group by starting inspectable parameters at $\bot$, accumulating
observations in fields/alias bodies and joining until stable. An
unobserved argument at an SCC-local occurrence contributes no observation.
After convergence phantom parameters default to $+$. Opaque and GADT
parameters start and stay invariant. This is an analysis for generalization,
not an assignment conversion $D\,s<:D\,t$.

Let $H$ be the union of the free variables of the types of expansive
bindings, intersected initially with the local variables; if a weak
candidate later becomes a compound type, track its variables. At the
final substitution, let $v(a,\bar u)$ be the sign joined over the
whole group. Define the eligible local variables as
\[
 L^+ = L\setminus
   \{a\in H : v(a,\bar u)\ne +\ \lor\ a\in\fv(P)\}.
\]
Here all sets in the formula use their substituted variables; eligible open
origins are handled separately from structural occurrences. When
generalization is disabled $L^+=\varnothing$. The \emph{monomorphism
restriction} is a second, distinct decision. The group is unrestricted
when generalization is enabled and all its members are functions; otherwise
it is restricted. Quantifying only eligible occurrences, for each $u_i$:
\[
\begin{array}{c|c|c}
 \text{group} & \text{scheme} & \text{outgoing predicates}\\ \hline
 \text{unrestricted} & \forall L^+.\ P\Rightarrow u_i & D\\
 \text{restricted}
    & \forall(L^+\setminus\fv(P)).\ \varnothing\Rightarrow u_i & D\cup P
\end{array}
\]
Only flexible variables are abstracted. All bindings of the group see the
same provisional context. Recursive occurrences are rewritten to carry the
solved type and needed dictionaries. Local weak variables can be solved by
later uses; top-level inferred schemes may contain neither free variables
nor predicates left for a caller to supply outside the scheme.

\begin{example}
An inferred numeric function $(\mathsf{fn}\ (x)\ (+\ x\ x))$ can keep
$\forall a.\mathsf{Num}\,a\Rightarrow a\arr a$; an implicit top-level
nonfunction initialized with $1$ defaults to $\Int$ in the usual library.
For a \emph{local} binding initialized with \code{(cell:new None)}, the provisional type
is $\mathsf{Cell}\,(\mathsf{Optional}\,a)$. $\mathsf{Cell}$ is opaque, so
$a$ is invariant and remains shared. Giving each lookup a fresh $a$
would let two consumers disagree about the type of the same mutable slot.
The same unannotated allocation at top level is rejected if its element
type remains unresolved; it does not introduce a top-level polymorphic cell.
By contrast an ordinary call returning $\List\,a$ with no predicate on
$a$ can generalize $a$ if it is local and there are only positive
occurrences in the group.
\end{example}

\begin{invariant}[Sharing of a blocked weak variable]
At an implicit binding boundary, suppose a remaining ordinary variable
$a$ of an expansive initializer is local and is either observed negatively
or invariantly in the group, or appears in a retained predicate. No scheme
created by the group makes fresh copies of that $a$ at lookup.
\end{invariant}
\begin{proof}
After substitutions the variable is in $H$. A negative occurrence joins
to $-$ or $\pm$; an invariant occurrence joins to $\pm$. Either excludes
$a$ from $L^+$, as does membership in $\fv(P)$. Neither row of the table
quantifies anything outside $L^+$. In the stored body its occurrences
therefore remain the same free identity, not a $g_i$.
$\operatorname{inst}$ freshens only abstracted positions (and distinct
origin metadata); any later ordinary equation constrains that same identity.
This is a sharing invariant, not a semantic proof that all positive
constructors or all foreign operations are safe.
\end{proof}

\paragraph{Declared bindings.}
For $\sigma$ freshen the declared binders, check the binding at the
instantiated body with its declared predicates available, solve its result
proofs, FDs and accessors, align determined arguments, and remove
entailed wanteds. Default and split the remainder. Requantify the
\emph{final} declared instance, excluding variables fixed to a structure,
identified with one another, or free in the outer environment. For explicit
binders use their declared order on the eligible survivors. The resulting
scheme must pass \src{ty-scheme=} against the substituted original; a
remaining retained predicate is a missing declared constraint. Thus an
annotation does not make an unconstrained expansive initializer more
general. Qualified declarations introduce dictionary parameters and do
not apply the same once-evaluated initializer restriction. The local
checker for a bare nested $\forall$ is different (next subsection).
Self-calls of a declared function instantiate its original scheme, not the
single instance being checked: polymorphic recursion needs this distinction.

Recursive SCCs of functions are allowed. Mixed mutually recursive functions
and values are rejected, as are recursive top-level values; there is a
special local value case for saturated, representation-eligible recursive
constructors with guarded components. The typed top-level/method recursion
check accounts for implicit dictionary parameters. These are initialization
conditions in addition to type equality.

\sources \src{define.lisp}: \src{infer-bindings-type},
\src{infer-impls-binding-type}, \src{infer-expl-binding-type},
\src{generalizable-type-variables}, \src{check-for-invalid-recursive-scc};
\src{context-reduction.lisp}: \src{reduce-context}, \src{split-context},
\src{candidates}, \src{default-builder-context}; \src{variance.lisp}:
\src{collect-tyvar-variances}; \src{define-type.lisp}:
\src{infer-type-definition-scc-variances}.

\subsection{Expected universals and expression ascriptions}\label{sec:rank}
Nested universals are checked, not discovered by guessing where to insert
$\forall$. When the substituted expectation is $\forall\bar a.\,P\Rightarrow u$,
intercept the ordinary check except at transparent body/block/ascription
nodes. The logical introduction rule is
\[
 \rrule{Forall-in}{\bar s\text{ fresh rigid}\quad
       G,P[\bar s/\bar a];\Gamma\vdash e\chk u[\bar s/\bar a]
       \quad\text{scope and value conditions}}
      {G;\Gamma\vdash e\chk\forall\bar a.\,P\Rightarrow u}.
\]
Specifically, snapshot the next fresh ID $m$ and outer bindings, scope the
names to $\bar s$, and check with a private result-proof scope. Resolve its
saved branch predicates before relaxing any skolem. Solve given and wanted
FDs, attempt accessor solving (an unresolved source involving a local
skolem is ambiguous), and align FDs from the givens. For every substitution
whose domain was allocated before $m$, inspect the substituted range; also
inspect substituted outer environment variables. Neither may contain a
local skolem. Discharge entailed predicates; after builder defaulting and
context reduction no remaining wanted may mention a local skolem.
If there are no given predicates, the expression is expansive and there
are skolems, every local skolem must be covariant or absent in the checked
body type. If there are givens, the elaborated expression instead accepts
dictionaries: attach them to an abstraction or insert a dictionary-binding
wrapper. Finally replace the local skolems by fresh flexible names inside
the elaborated object; the exported type is the substituted \emph{original}
$\forall$. This administrative relaxation is not an equation on caller types.

\begin{example}
To check \code{(fn (z) z)} at $\forall a.\,a\arr a$, use a fresh $s$;
the parameter and result both have type $s$ and no prior unknown must
equal it. Conversely, if $x:b$ is a captured monomorphic value, checking
\code{(fn (\_) x)} at that type would solve the old $b$ to $s$ or require
an incompatible rigid equality. The former fails the watermark test.
The use of a variable bound to the successful polymorphic parameter can
instantiate it independently at $\Int$ and $\Str$.
\end{example}

\paragraph{The \code{the} form is not an unchecked assertion.}
Parse its annotation with the current scoped variables and branch
refinements. A predicate context requires explicit quantification. If
the declaration is explicitly quantified and either has predicates or is
checked at a $\forall$ expectation, use universal introduction on the
operand and compare the resulting
polytype (or instantiate it if an ordinary type is expected). In the
other case let $d$ be the parsed declared body. Check compound operands at
$d$ so e.g.\ a match sees its expected result. Simple leaves without
nested $\forall$ are inferred at a fresh result-capable expectation.
Let $u$ be the resulting \emph{substituted} operand type. Require
$u\eqs d$ and also $u\oneway d$ on the saved operands (not their common
unified image), then unify $u$ with the enclosing expectation. The
one-way check prevents certifying an annotation just because it could be
made \emph{less} general to fit. There is no recursive subsumption rule
that moves binders through a function, reorders visible arguments, or
delays an effect to manufacture a more polymorphic value.

\sources \src{define.lisp}: \src{check-forall-expression-type},
\src{parse-the-type-annotation}, \src{node-the} method of
\src{infer-expression-type}, \src{wrap-expression-with-given-predicates};
\src{tc-env.lisp}: \src{tc-env-extend-type-variable-scope}.

\section{Constructors, patterns and local equality}\label{sec:patterns}

\subsection{Introduction versus opening}
An ordinary constructor $K$ of $D\,\bar a$ with fields $\bar t$ has the
scheme $\forall\bar a.\,\bar t\arr D\,\bar a$ (the bare result type for a
nullary constructor). A struct is nominal, with a constructor and named
fields. With the GADT spelling each constructor instead has an independently
scoped signature
\[
        K:\forall\bar b.\,\bar Q(\bar b)\Rightarrow
          \bar t(\bar b)\arr D\,\bar u(\bar b).
                                                        \tag{K}
\]
Some indices may be concrete, repeated, or structured; some binders may
occur only in fields or predicates. The signature must return exactly one
instance of the defining nominal type, of kind $*$; fields are positional.
A nullary constructor uses the bare result type, not a nullary function
returning it, and remains source-nullary even when $\bar Q$ is nonempty.
The implementation rejects mixed ordinary/signed constructor lists; each
constructor signature is checked under the standard kind, scope, ambiguity,
and functional-dependency rules, and a polymorphic field can itself carry a
qualified $\forall$. At construction the caller instantiates all binders
together, supplies dictionaries for $\bar Q(\bar b)$, checks the visible
fields $\bar t(\bar b)$, and obtains $D\,\bar u(\bar b)$ under the same
instantiation. There is no mystery type chosen by the consumer. Physical
dictionary slots precede the visible fields in the constructor's runtime
representation while remaining hidden from source arity and pattern syntax;
unconstrained constructors keep their layout unchanged. Explicit
representations that cannot carry hidden dictionary slots
(\code{:enum}, \code{:transparent}, and \code{:native}) reject constrained
constructors.

For elimination write
\[
       \Delta;E\vdash p:t\ \rightsquigarrow\ (\Delta_p,E_p,G_p,\Gamma_p).
\]
This is the \emph{meaning} of a pattern assuming that its runtime test
succeeded. A wildcard adds nothing; a variable binds the input type; an
as-pattern binds the whole value as well as its subpattern and retains its
subpattern's evidence. A literal tests a value of the input type. For
$K(\bar p)$ against $D\,\bar v$, open (K) at fresh local names $\bar b$
and add $\bar v=\bar u(\bar b)$ to $E_p$ and the stored dictionaries
$\bar Q(\bar b)$ to $G_p$ for the dependent subpatterns and continuation
only; then check the field patterns at $\bar t(\bar b)$ in this same row.
For a branch with outer result $r$ the meaning is
\[
 \rrule{Arm}{\Delta;E\vdash p:t\rightsquigarrow(\Delta_p,E_p,G_p,\Gamma_p)
          \quad \Delta,\Delta_p;E,E_p;G,G_p;\Gamma,\Gamma_p\vdash e:r}
      {\Delta;E;G;\Gamma\vdash (p\mapsto e):t\Rightarrow r}.
\]
Names and dictionaries introduced by opening are bound in this premise, not
in $t,r$ or a sibling premise; two independently matched values open
distinct local names. An obligation about such a name cannot be moved
outside as an implicit universally quantified predicate. Constructor return
indices affect compatible type instances, not the identity of the enclosing
type constructor. E.g.\ for $K:\forall h.\,h\arr D\,(\List\,h)$, matching
$D\,a$ gives $a=\List\,h$. One can return something of type $\List\,h$
under the \emph{outer} name $a$; one cannot return an $h$ under an
arbitrary caller-chosen type.

\subsection{The two environments of a GADT match}
For a match without GADT patterns infer the scrutinee at a fresh ordinary variable,
unify its type against every pattern and check every arm at the same
expected result. If \emph{any} pattern contains a signed constructor
(also inside an ordinary constructor or as-pattern), use the following
branch-local algorithm for the whole match.
\begin{enumerate}
\item Solve available scrutinee accessors. Compute a common base substitution
$S_0$ by collecting input \emph{shape} requirements across rows. E.g.\
knowing only $x:a$, a pattern may require $x:D\,b$ for fresh $b$;
it does not demand of all rows that $b=\Int$ just because this constructor
returns $D\,\Int$. This preliminary walk uses a local solution for nested
witnesses while committing required nominal shapes.
\item For each arm start the pattern from the same $S_0$ in a copied term
\emph{and} scoped-type environment. Snapshot the next variable ID $m$.
Freeze flexible variables of the substituted scrutinee as fresh skolems.
When a GADT result and input have matching applied heads, use
\src{unify-given}: temporarily open permitted rigid variables on the
two sides, unify the results, and close each surviving abstract flexible
representative back to a skolem. Otherwise ordinary unification must
justify the head. Only a constructor test has this authority; a term,
a literal test and a wildcard do not.
\item Open each constructor at the exact field count. To describe
\src{open-constructor-pattern}, instantiate the constructor's fields,
result and predicates (together with their superclass evidence) at the same
fresh binder variables. Let $V$ be the variables of the incoming
substituted input, and $P$ the substitution after the constructor's result
has been unified with it. Of the nonrigid variables in the substituted
constructor type and the superclass closure of its predicates, seal those
\emph{not equal to an element of} $[Pv:v\in V]$ as fresh existential
skolems $H$. This is not set subtraction of $\fv(PV)$: an unnamed component
of a newly exposed $\List\,h$ or a hidden superclass witness may still
need sealing. Store the constructor binder positions that open to sealed
names on this occurrence, and record the refined constructor predicates on
the typed constructor pattern. Children may use permitted row evidence,
including such witnesses from enclosing constructors.
\item A row uses one queue for all columns, also for parameters. Choose a
GADT-containing pending pattern with a known nominal input head if possible,
otherwise the first witness, otherwise the first item. Replace a visited
constructor by its fields in place. This allows a tag in one field to identify
the type at which a sibling literal is to be tested. Typed field and parameter
order remains source order. Rename residual index representatives back into
the outer vocabulary when possible. Call the resulting pattern substitution
$P$ and retain $E_p=\{v=Pv : v\in\fv(S_0t),\,v\ne Pv\}$ plus enclosing
evidence. Normalize the branch environment, its scoped annotations, the
recovered constructor predicates $G_p$, and enclosing dictionaries by $P$
(recording $H$ as unnamed fixed scope variables) before checking the body.
\item Check the body at the locally refined expected result, with the
following exception. Use the contextual return target when applicable.
If the resulting expectation under $S_0$ is a variable not occurring in
the substituted scrutinee, and local or enclosing index evidence is
available, check at a fresh flexible copy and \emph{project} the inferred
type outward as described below.
Let $T$ be the final body substitution. Export as body requirements only
equations $v\mapsto u$ in $T$ with $\operatorname{id}(v)<m$ and
$v\notin\dom P$. Re-unify them with the accumulating outward substitution;
do not export the domains of $P$ merely because later composition changed
their ranges. A sibling's own pattern starts from $S_0$, not $P$ or $T$.
\end{enumerate}

\paragraph{Choosing an outward result.}
Given a locally normalized result type $u$, for an occurrence of a subtree
$w$ collect its aliases $\{a:a=w'$ is in the saved evidence and
$Tw'=w\}$. Recursively project arguments, function slots and pack components
to obtain a structural spelling $w'_{\rm str}$. With no aliases, use that
spelling (escape checks still apply). Otherwise the candidates are the
aliases and, unless it mentions a hidden variable, $w'_{\rm str}$. With
several candidates record a fresh \emph{monomorphic} choice together with
the ordered options (structural first when permitted). At boundary solving,
accept an already equal option, otherwise prefer a unifiable option shared
by later choices for the same wanted type, then try the first unifiable
option, failing on incompatible types. This is a finite choice procedure,
not general anti-unification. A local binding may not quantify an undecided
choice or the variables it depends on. A nested arm with hidden variables
finishes nested choices before its escape check.

\begin{lemma}[Projection respects the evidence]\label{lem:project}
On closed-interface first-order types, suppose a valuation satisfies the
arm's saved index equations under its body substitution and a solution of
all its recorded choices (each choice equals one of its candidates).
Interpreting the projected result under that valuation gives the same type
as interpreting the local result.
\end{lemma}
\begin{proof}
Induct bottom-up on the local type. A structural spelling rebuilds the same
head from equal children by the induction hypothesis. Every alias was
added only by equality with that very subtree after the body substitution;
the corresponding saved equation makes the alias denote the same tree.
Omitting an option that mentions a hidden variable does not add any new
interpretation. Replacing the remaining options by a variable equated
to one of them therefore preserves equality. The assertion is local:
it does not allow substituting this arm's $a=\Int$ into another arm.
\end{proof}

\paragraph{Predicate and existential discharge.}
Suppose an arm checks a use requiring $C\,a$ and the enclosing or
constructor-recovered dictionary has type $C\,b$. A constructor proving
$a=b$ can make these the \emph{same} predicate locally; merely being of
class $C$ cannot. Before discharging branch predicates, fresh flexible
branch unknowns determined by the arm's given predicates via functional
dependencies are solved by one-way matching against those givens.
Obligations mentioning hidden types are solved in this scope. Repeatedly
solve fields whose receiver mentions a hidden type. For hidden predicates
try supplied and constructor-recovered dictionaries and their superclasses,
then stably selected instance prerequisites; keep unproved predicates at a
repeated goal or instance-path depth 64. Nonhidden prerequisites not
already entailed by the arm's given predicates can still go outward. After
substitution, \src{check-pattern-scope} rejects a hidden variable in the
outward result, supplied extra parameter types being checked, an old
substitution range, an unsolved predicate, or either endpoint of a pending
accessor.

Let $A\,k\,x\,y$ declare $k\to x\,y$. Suppose supplied or
constructor-recovered dictionaries $F\,q\,x$ and $S\,k\,y$ expose
superclass fields $A\,q\,x\,h_1$ and $A\,k\,h_2\,y$. Inside the success arm
of $\mathsf{Same}\,q\,k$ the determinants coincide. On an indexed match or
positional parameter edge, \src{refine-given-superclasses} uses the
available predicates' stored superclass fields at the local types,
materializing missing projections and marking the available dictionaries as
\src{pattern-given} for use on this edge. From these predicates and their
actual fields, \src{pattern-given-equations} computes $E_G$ by pairwise FD
improvement to a fixed point. Equal determinants justify equality of each
declared output, including two distinct names fixed by the caller.
These are equations for \emph{comparing evidence}, not assignments in the
type inference substitution. If the output equations conflict, no FD
conversion is used. A wanted predicate and a global instance are never
premises of this calculation.

For example, $h_1=y$ and $h_2=x$ after the key match. Now add supplied
$F\,k\,z$ and $S\,q\,w$. Prior to the match the two pairs of dictionaries
may already have fields $A\,q\,x\,w$ and $A\,k\,z\,y$. After $q=k$
the FD proves $x=z$ and $w=y$; either actual field agrees with the
request $A\,k\,x\,y$ under $E_G$. Binding inference need not have kept
that literal spelling at any field. The checker and translator compare
both sides under the equations of their respective givens; the translator
prefers an exact path, otherwise projects a field of equal type from the
original dictionary. Under the hypotheses of Lemma~\ref{lem:fd} each
FD comparison preserves every ground solution, hence representative
choice does not affect ground equality. This does not assert uniqueness
of dictionary values or completeness of inference. Retained result
obligations use their saved givens (including nonhidden constructor-local
givens of that arm) and index equations. No conversion is passed to
instance selection, a sibling arm, or the caller's solver. One partial
dictionary alone does not let the caller choose the other output.

An overloaded literal can occur \emph{before} the key or dictionary witness
in source order. Its operand type and even the head of its captured
dictionary may already be known while a constructor-bound dictionary or the
type of its numeric superclass field is not. The row scheduler defers any
literal test that references a not-yet-matched constructor's dictionary
binding and checks the selected accessor path from the active dictionaries
using only the equations established by the current pattern prefix. Thus it
need not run an effectful comparison on a path where its required
dictionary or key equality has not been established.

When the result itself is being chosen, a predicate with fresh variables
may depend on a not-yet-known result. At an explicit binding or an expected
$\forall$ the checker can retain $(p,G_p,E_p,H)$ as an \emph{obligation},
not as an equation or a dictionary. It does so for nonhidden predicates
under its retention policy when their fresh variables occur in the result,
or when a supplied predicate can be matched by binding only fresh result
variables. Independently synthesized functions do not inherit a contextual
result merely by being lexically nested. Once the result substitution $S$
is known, transport saved equations on \emph{both} sides
($a=u$ becomes $Sa=Su$), normalize wanted and givens and test entailment.
Failed transport on a contradictory branch leaves compatibility/generality
to the ordinary checks. An unproved predicate is returned for ordinary
checking (or rejected for escape); it is not deemed proved by retention.
At a rank boundary only unresolved obligations independent of the local
skolems, with a retaining parent, can be forwarded; the implementation
filters out givens mentioning those skolems. It finalizes the typed arms
using their local equations and stable branch identities before relaxing
skolems and abstracting dictionaries. Equality evidence need not survive
as a runtime object; dictionary evidence does.

\begin{example}[Same tag test, two different result indices]
The expected output type can be given the scrutinee's index:
\begin{verbatim}
(define-type (Atom :a)
  (declare AnInt (Integer -> Atom Integer))
  (declare AString (String -> Atom String)))
(declare unpack (forall (:a) (Atom :a -> :a)))
(define (unpack term)
  (match term ((AnInt x) x) ((AString y) y)))
\end{verbatim}
The shared input is $\mathsf{Atom}\,a$, the expected result $a$. The
first arm has $a=\Int,\ x:\Int$, so $x:a$ by its equation; the second
has $a=\Str,\ y:\Str$. Discharging these two implications gives the
declared result without solving $\Int\eqs\Str$. A wildcard fallback would
get \emph{neither} positive equation merely from previous patterns failing.
Conversely, against a statically known $\mathsf{Atom}\,\Int$, an explicitly
written \code{AString} arm fails pattern compatibility; there is no rule
allowing arbitrary code because its premises are contradictory.
\end{example}

Sequential GADT pattern bindings give the continuation the arm scope by
rewriting it as a one-arm match. For monadic \code{<-}, the action stays
in the prefix; matching is in the callback, not in a computation that
repeats the action. Parameter patterns similarly check their common
input shapes, body, keyword defaults and result with local evidence;
catch/resumption pattern scopes have their own checks. In a separate
analysis pass the pattern matrix specializes candidate constructor results
against a column type and propagates equations to remaining columns.
It checks compatibility again after later uses may have specialized
the scrutinee. A non-exhaustive \code{match} expression produces a
\emph{warning}; a refutable function parameter row or a refutable
\code{let} binding pattern is an error. Do not infer totality from the
success of expression typing.

\sources \src{define-type.lisp}: \src{check-constructor-signature},
\src{infer-define-type-scc-kinds}; \src{unify.lisp}: \src{unify-pattern-result},
\src{unify-given}; \src{define.lisp}: \src{constrain-pattern-inputs},
\src{open-constructor-pattern}, \src{infer-pattern-row}, \src{branch-body-equations},
\src{project-branch-result}, \src{solve-branch-result-choices},
\src{resolve-branch-result-predicates}, \src{discharge-hidden-obligations},
\src{check-pattern-scope}; \src{src/analysis/analysis.lisp}:
\src{analyze-translation-unit}; \src{src/analysis/pattern-exhaustiveness.lisp}:
\src{possible-constructor-pattern}, \src{find-non-matching-value}.
The input-shape and pattern-row walks both use
\src{visit-pattern-queue}; matrix specialization and defaulting share
\src{refine-column-types} to transport the remaining column types, fields
and obligations.

\section{Control, effects and the host}\label{sec:effects}
There is no effect row in $\Fn$. A type records what a computation may
return \emph{normally}, not that it is pure, terminates, cannot signal a
condition or is being run in the dynamic extent of a handler.

\paragraph{Local choice and iteration.}
\code{if} and each \code{cond} clause require Boolean tests and check result
arms at the common expectation. \code{and}/\code{or} take and return
Booleans. \code{when}/\code{unless} check a Boolean test, discard any body
results and return $\Void$. In \code{for}/\code{for*} mutable loop
variables have \emph{one} instance for initializer and step (generalization
disabled even for declarations). A while/until test has type $\Bool$, a
repeat count $\mathsf{UFix}$, and the body checks at $\Void$. The loop's
result comes from the optional returns form or is $\Void$; \code{break}
and \code{continue} check as $\Void$, not as expressions of every type.
\code{rec} binds and immediately invokes a local tail-recursive function.
The independent-initializer path seeds monomorphic parameter types from
the initial call; dependent initializers use the nested-let lowering.
Only direct self-calls in allowed tail positions are admitted.

\paragraph{Named exits.}
\src{resolve-control-flow} rewrites surface \code{return} to a named exit
from its lexically enclosing return scope, rejecting it in disallowed
contexts. A block introduces a shared result-capable target $B(\ell)$
unified with its expectation and the normal tail. The meaning rule is
\[
 \rrule{Exit}{B(\ell)=u \quad \Delta;E;B;\Gamma\vdash e:u}
       {\Delta;E;B;\Gamma\vdash \mathsf{returnFrom}(\ell,e):t}.
\]
There are two result types here: the payload's destination $u$ and the
surrounding expression's expectation $t$, to which the exit does \emph{not} return.
Ordinarily the checker infers payloads at fresh result-capable types,
records the first, then checks later returns and normal completion against
it. At a target with local GADT evidence, check the payload at the locally
substituted target directly; do not record a branch-only spelling as the
type of other exits. A $\forall$ target also supplies the operand
expectation. Fresh-result-only hints may carry the destination through
intervening function results. A new lambda introduces its own return
scope; returns in parameter defaults share its body block.

\paragraph{Classes of effects.}
In \code{do}, choose $m:*\arr *$, demand $\mathsf{Monad}\,m$, check each
action at $m\,a_i$, the bound pattern at $a_i$, and the last expression at
$m\,b$ for the result $m\,b$. Ordinary local bindings are still possible
inside the sequence. This is an overloaded sequencing interface, not an
effect discipline for Coalton expressions. Collection and association
builders check homogeneous items (respectively keys and values), then
typecheck the generated builder-protocol calls; comprehension clauses
likewise lower to library operations. Dynamic rebinding checks an
initializer in the \emph{old} environment, applies generalization/defaulting
and requires the resulting scheme to equal the variable's existing scheme.

\code{throw} accepts a known exception type, or a variable with an
$\mathsf{Exception}$ predicate; \code{resume-to} analogously uses
$\mathsf{Resumption}$. Both have the local expected result, since they
transfer control. \code{catch}/\code{handle} check their guarded expression
and every handler body at one result (including a pack); their patterns
are restricted to admissible exception tests, optionally typed, or
wildcards. \code{resumable} requires a \emph{single} result for its guarded
expression and branches, and resumption-constructor patterns. These
checks do not ensure that a dynamically applicable handler exists.
\code{protect} preserves the guarded expression's result sequence and
discards the cleanup result. \code{type-of} infers its operand, solves
its constraints/accessors and returns a quoted $\mathsf{TypeScheme}$;
it does not evaluate that operand to determine its type.

\paragraph{Foreign terms and representation.}
A \code{lisp} form declares zero, one or several outputs. Parse these
types in the lexical type-variable environment, form $\pack(\bar t)$ and
unify with the context. Look up the explicit captures; a fresh constrained
capture may specialize to an unambiguous supplied predicate, but may not
solve an old lexical variable to do so. A still-fresh ambiguous capture
predicate is an error, not a defaulting request. The checker does
\emph{not} derive a typing for the host forms themselves. The host must
deliver the claimed representation and arity, respect captured values and
obey representation invariants. For example, a \code{lisp} form claiming
$\Int$ can contain host code producing a string; the declared output
alone cannot prove otherwise. Generated checks may signal errors, and
\code{unsafe} changes host safety settings without disabling Coalton's
static checking. A \code{repr} declaration, optimizer specialization or
host-side mutation is not validated by a source substitution lemma.

In particular, \code{specialize} checks that the replacement value's scheme
equals the requested unqualified specialization, is not scheme-equal to
the original, and that the original's fresh body type one-way matches the
replacement's. This is a type test, not observational equivalence of the
two definitions. In development mode redefinition checks guard, among
other things, changes to the representation/constructor evidence of a
GADT. The metatheory below fixes the signature throughout a run.

\sources \src{control-flow.lisp}: \src{resolve-control-flow};
\src{define.lisp}: \src{infer-rec-node-type}, \src{infer-dynamic-binding-type},
\src{reflect-expression-type-scheme}, \src{specialize-lisp-capture},
\src{infer-expression-type} methods for the named nodes;
\src{specialize.lisp}: \src{process-specialize}; \src{lisp-type.lisp}:
\src{lisp-type}; \src{define-type.lisp}: \src{validate-gadt-redefinitions}.

\section{A type safety result with an explicit boundary}\label{sec:safety}

\subsection{The indexed, single-result core}
Fix a signature of free nominal datatypes, distinct atomic base types for
primitive literals, and unconstrained constructor signatures (K)
($\bar Q=\varnothing$) whose parameters bind all the variables in their
field and result types. All type parameters here have kind $*$. Types are
first-order trees $a,\ B,\ D\,\bar t,\ (\bar t\arr t)$; in this notation
nominal applications are saturated, arrows are closed and the result is
always one value. Recursive nominal declarations are fine; there is no
term-level recursion, recursive let, cyclic value construction or
fixed-point primitive in this core. Terms are variables, fixed-type
constants, uncurried lambdas with variable parameters, exactly saturated
calls, saturated constructor introductions, and cases with distinct
constructor tags and variable fields. Write a nullary constructor as
$K()$. Parameters in a binding list are distinct. There is no type-level
computation except finite tree equality. No overloaded literals, class
methods, constructor-local constraints, keyword arguments, generalization
of local bindings, nonlocal control, host code, mutation, or optimizer
rewrites occur in this result. The omitted constructs keep the checking
rules above; they do not acquire a semantics by erasure.

For a kinded variable context $\Delta$, let $\mathrm{Types}(\Delta)$
denote the well-kinded core types with free variables in $\Delta$.
Write $\rho:\Delta$ for an assignment of closed types of the required kinds. Define
$\Delta;E\models t=u$ to mean that $\rho t=\rho u$ for every such
$\rho$ satisfying all equations in $E$. Equality here is equality of free
trees, not equality of host representations. Define
$\Delta;E;\Gamma\vdash_0 e:t$ by the following rules (suppressing the
fixed $\Sigma$). Variables have their types in $\Gamma$ and constants
their fixed base types. Unlike the checking rules of
Section~\ref{sec:terms}, these rules have no substitution state:
\[
 \rrule{Fn}{\Delta;E;\Gamma,\bar x:\bar t\vdash_0 e:u}
           {\Delta;E;\Gamma\vdash_0 \lambda\bar x.e:\bar t\arr u}
\]
\[
 \rrule{Call}{\Delta;E;\Gamma\vdash_0 f:\bar t\arr u\quad
              (\Delta;E;\Gamma\vdash_0 e_i:t_i)_i}
             {\Delta;E;\Gamma\vdash_0 f(\bar e):u}.
\]
The two lists in a call have the same length, including length zero.
\[
 \rrule{Conv}{\Delta;E;\Gamma\vdash_0 e:t\quad\Delta;E\models t=u}
             {\Delta;E;\Gamma\vdash_0 e:u}
\]
\[
 \rrule{Make}{K:\forall\bar b.\,\bar t\arr D\,\bar u
       \quad\eta:\bar b\to\mathrm{Types}(\Delta)
       \quad\Delta;E;\Gamma\vdash_0\bar e:\eta\bar t}
       {\Delta;E;\Gamma\vdash_0 K(\bar e):D\,(\eta\bar u)}.
\]
Parentheses in $D\,(\eta\bar u)$ delimit its vector of indices. A case
at $r$ requires a scrutinee at $D\,\bar v$ and, for each written
alternative $K_i(\bar x_i)\mapsto e_i$, the premise
\[
 \Delta,\bar b_i;\
 E,\bar v=\bar u_i;\
 \Gamma,\bar x_i:\bar t_i\ \vdash_0 e_i:r,                  \tag{Case}
\]
where every written $K_i$ belongs to $D$, has the displayed field arity, and
$\bar b_i$ freshen \emph{its} signature $K_i:\forall\bar b_i.\,
\bar t_i\arr D\,\bar u_i$. Thus this premise cannot smuggle an equation
from a constructor of a different datatype into the body. Neither $\Gamma$
nor $r$ mentions these new binders. Impose the following sufficient
\emph{coverage} premise: for every $\rho:\Delta$ satisfying $E$, every
constructor whose result can be instantiated to $D\,(\rho\bar v)$
has a written alternative. This is a mathematical premise, not an appeal
to the absence of compiler warnings. It may conservatively require a
constructor even when one of its fields is uninhabited.

\begin{proposition}[Monomorphic checker fragment]\label{prop:mono}
Take a successful check of the core syntax with monomorphic term entries (apart from
constructor instantiation), no remaining predicates/accessors, only the
ordinary equations of Theorem~\ref{thm:mgu}, and no overloading, hidden
keyword interfaces or result choices. For a case require that its
scrutinee indices and expected result are already ground when the patterns
are checked, and that the stated coverage premise holds. Saturated
constructor syntax here abbreviates its ordinary lookup and exact call.
Apply the solved substitution and complete any remaining flexible unknowns
with well-kinded closed core types consistent with it. Rigid witnesses
introduced in case arms stay bound there. The erased term at these types
has a $\vdash_0$ derivation.
\end{proposition}
\begin{proof}
Induct on the successful check. For a leaf take its substituted environment type or fixed literal
type. For a lambda extend by the solved positional types and use the
derivation for the body. For a call use the derivations for its operator
and operands; the solved slot equalities give the premises of \textsf{Call}.
For a constructor the instantiation and those same equalities give
\textsf{Make}. In each case Theorem~\ref{thm:mgu} justifies the equation
step, including a successful operator-result hint; no operand proof
depends on the traversal order.

For a flat arm against $D\,\bar v$ with $\bar v$ ground there is
no outer index to solve by freezing or projection. Unifying a fresh
constructor result with $D\,\bar v$ fixes every constructor binder
appearing in that result to the corresponding ground subtree; repeated
occurrences must agree. The remaining field-only binders are unconstrained
by the result and are sealed as distinct skolems. The body's recursive
derivation may use the fixed subtrees and these arbitrary representatives.
Under $\bar v=\bar u_i$ the fixed binder positions equal those subtrees
by injectivity, so \textsf{Conv} justifies their uses. Renaming each
sealed representative to its fresh $\bar b_i$ binder gives the other
field types in (Case). The ground $r$ is identical outside. Combining
the arm derivations with the scrutinee and coverage gives the case rule.
\end{proof}
This induction on the ordinary checker cases does not yet justify an
elimination whose scrutinee has a rigid, \emph{unknown} index.

\subsection{From a rigid index to a core derivation}\label{sec:elim}
Consider a separate class of inputs for which we can follow that path all
the way through the checker. Let a fresh unary $D$ have finitely many
constructors of the following forms. Take at least one of each so that
the checker encounters both kinds of result projection:
\[
  I_i:B_i\arr D\,B_i,\qquad
  K_j:\forall b.\ b\arr D\,(F_j\,b).
\]
Here $B_i$ are ground base types, and each $F_j$ is a unary ordinary
datatype with a constructor $A_j:\forall b.\ b\arr F_j\,b$. Require fresh
term binders, all tags exactly once, and no other fields, constraints,
captures, annotations or operations in the following template. A declared
nullary function returns a lambda at $\forall s.\ D\,s\arr s$. Its body
matches its parameter $p$, returning $x$ in $I_i(x)$ arms and $A_j(x)$ in
$K_j(x)$ arms. It binds that match to a local monomorphic $z$ and returns
$z$. One complete input, with declared arities and no foreign operations:
\begin{verbatim}
(define-type (Envelope :a) (Enclose :a))
(define-type (Parcel :a)
  (declare Counted (Integer -> Parcel Integer))
  (declare Wrapped (:b -> Parcel (Envelope :b))))
(declare make-unpacker
  (Void -> (forall (:a) (Parcel :a -> :a))))
(define (make-unpacker)
  (fn (p)
    (let result = (match p
      ((Counted n) n)
      ((Wrapped x) (Enclose x))))
    result))
\end{verbatim}
Putting the universal \emph{inside} the declaration matters. A prenex
\code{declare} freshens its variables and rechecks generality afterwards;
the result position here actually calls \src{check-forall-expression-type}
and allocates a rigid skolem. Naming the match result matters too: its
immediate expectation is fresh, not simply the known answer $s$.

\paragraph{What the checker projects.}
In a successful check of this template the match projects to the declared
index without exporting an equation on that index. The universal check gives
$p:D\,s$ for fresh rigid $s$ and demands body result $s$. The binder for
$z$ makes the match synthesize at a fresh result-capable $r$. The common
shape pass already sees $D\,s$; freezing has nothing to do to $s$, and
the local pattern checks start from this same shape in each arm. Since
$r$ is distinct from the input index, local evidence triggers projection.
\begin{itemize}
\item For $I_i$, \src{unify-given} temporarily opens $s$ to a flexible and
unifies it with $B_i$. The resulting local equation is $s=B_i$;
the bound $x$ has type $B_i$. There is no hidden variable. Projecting $x$'s
type yields a fresh choice $c_i$ with options $[B_i,s]$.
\item For $K_j$, opening the signature introduces a fresh \emph{flexible}
$b$. Equating $D\,s$ with $D\,(F_j\,b)$ under the permitted opening of
$s$ exposes $s=F_j\,b$. Neither $s$ nor its whole image $F_j\,b$ names
the component $b$; \src{open-constructor-pattern} seals $b$ as a fresh
rigid $h_j$. The field thus has type $h_j$ and local evidence is
$s=F_j\,h_j$. Instantiating $A_j$ at a new flexible and checking the
field argument sets that flexible to $h_j$, giving $A_j(x):F_j\,h_j$.
For the whole result the structural spelling contains the hidden $h_j$;
the only alias is $s$. Projection therefore returns $s$ itself, not a
constraint that could bind the caller to $F_j\,h_j$.
\end{itemize}
No other syntax here produces an outward equation on $s$. In particular
the equation whose domain is $s$ is in the pattern substitution, so
\src{branch-body-equations} does not export it. Merging the projected
results equates $r$, each $c_i$, and $s$; the $K_j$ arms supply the latter
spelling directly. Ordinary variable unification can direct a flexible
toward $s$ but cannot bind $s$. Each $I_i$ choice now already equals one
of its options, so no $s=B_i$ is exported by choice resolution either.
The final occurrence of $z$ checks at $s$. All $h_j$ stay local, all
uses of $s$ outside the arms belong to the universal's check, and the
unconstrained lambda is nonexpansive. The universal is closed again on
exit. In the concrete input above the two projection calls respectively
produce a flexible choice and the rigid alias; their types are not both
guessed from the last use of \code{result}.

\begin{theorem}[Core typing of the eliminator template]\label{thm:elim}
In a fixed signature as above, let $e_\lambda$ be the template's returned
lambda, forgetting types and unused implicit return blocks. Encode the
strict binding $\mathsf{let}\ z=M\ \mathsf{in}\ z$ by $(\lambda z.z)(M)$.
Then
\[
  s:*;\varnothing;\varnothing\vdash_0 e_\lambda:D\,s\arr s.
\]
For every closed core type $t$, the same lambda has type $D\,t\arr t$.
\end{theorem}
\begin{proof}
Construct the derivation from the template and its constructor signatures,
not from checker success. Work in $\Delta=\{s:*\}$ with $p:D\,s$. In an
$I_i$ arm assume $s=B_i$ as in (Case). Lookup gives $x:B_i$; \textsf{Conv}
gives $x:s$. In a $K_j$ arm take a fresh $b_j:*$ with $s=F_j\,b_j$.
Lookup gives $x:b_j$ and \textsf{Make} with the explicit instantiation
$[b\mapsto b_j]$ gives $A_j(x):F_j\,b_j$. \textsf{Conv} gives $A_j(x):s$.
This conversion expresses exactly the outward projection just traced;
it does not promote either branch equation into $\Delta$. All tags of
$D$ are present, so (Case) gives $M:s$ without equations. Type the identity
at $s\arr s$, apply it to $M$, then abstract $p$ to obtain the displayed
judgment. For any closed $t$, use the type-substitution part of
Lemma~\ref{lem:substitution} with $[s\mapsto t]$: the outer equation
context is empty, so the substitution has no further premise to solve.
At runtime a matching constructor value supplies its own witnesses for
the branch equations. This specializes a \emph{derivation}; rerunning
source inference at a narrower scrutinee type might reject a now
incompatible written arm.
\end{proof}
Nested patterns, dependent dictionary obligations, nonlocal exits and
arbitrary result-choice search are not covered by this correspondence.

\subsection{Evaluation and proof}
Constants and lambdas are values; $K(\bar v)$ is a value when all fields
are values. Write $\longmapsto_h$ for the two head reductions:
\[
\begin{aligned}
 (\lambda\bar x.e)(\bar v)&\longmapsto_h e[\bar v/\bar x],\\
 \mathsf{case}\ K_i(\bar v)\ \mathsf{of}\ \{\ldots,
       K_i(\bar x)\mapsto e_i,\ldots\}
       &\longmapsto_h e_i[\bar v/\bar x].
\end{aligned}
\]
The evaluation contexts $\mathcal J$ choose the next position from left
to right; $\bar v$ is an already evaluated prefix and $\mathcal A$ is a
list of alternatives:
\[
\begin{aligned}
 \mathcal J ::={}&[\,]\mid \mathcal J(\bar e)\mid v(\bar v,\mathcal J,\bar e)\\
                 &{}\mid K(\bar v,\mathcal J,\bar e)
                   \mid \mathsf{case}\ \mathcal J\ \mathsf{of}\ \mathcal A .
\end{aligned}
\]
One evaluation step is $\mathcal J[e]\longmapsto\mathcal J[e']$ whenever
$e\longmapsto_h e'$. In particular no step enters a lambda or an untaken
arm. Term substitution is simultaneous and capture avoiding. Type witnesses
belong to derivations, not to the runtime test.

\begin{lemma}[Weakening]\label{lem:weakening}
Let $\Delta;E;\Gamma\vdash_0 e:t$. Extend $\Delta$ and $\Gamma$ with
fresh bindings, preserving the old kinds and types, and strengthen $E$
to well-formed $E'$ entailing each old equation in the extended type
context $\Delta'$. Then $\Delta';E';\Gamma'\vdash_0 e:t$.
\end{lemma}
\begin{proof}
Induct on the derivation. Old lookups and constant types are unchanged.
For \textsf{Conv}, an assignment to $\Delta'$ satisfying $E'$ restricts
to one satisfying $E$, so it preserves the required equality. All type
instantiations in \textsf{Make} are still well formed; rebuild the rule
on the weakened children, as for a call. At a lambda freshen its bound
term names and extend both term contexts by the same argument types.
At a case, any tag possible for a new valuation was possible for its
restriction, so coverage survives. Freshen each arm's type and term
binders away from the added bindings, keep the same field types and add
the same branch equations on each side. The extended new equations entail
the extended old ones, so induction gives the required arm derivation.
\end{proof}

\begin{lemma}[Substituting types and values]\label{lem:substitution}
If $\Delta;E;\Gamma\vdash_0 e:t$ and a type assignment
$\theta:\Delta\to\mathrm{Types}(\Delta')$ satisfies
$\Delta';E'\models\theta E$ (each substituted equation), then
$\Delta';E';\theta\Gamma\vdash_0 e:\theta t$.
If $\Delta;E;\Gamma,\bar x:\bar t\vdash_0 e:u$ and
$\Delta;E;\Gamma\vdash_0\bar v:\bar t$, then
$\Delta;E;\Gamma\vdash_0 e[\bar v/\bar x]:u$.
\end{lemma}
\begin{proof}
For the first assertion induct on the typing derivation. Lookup,
constants, lambda and call are structural. In \textsf{Make} use
$\theta\circ\eta$; in \textsf{Conv} any $\rho'$ satisfying $E'$ gives
$\rho'\circ\theta$ satisfying $E$. In (Case) freshen the local binders
$\bar b_i$ away from $\theta$ and extend it by the identity on them.
Only $\bar b_i$ can occur free in the signature's $\bar t_i,\bar u_i$,
so those field and index types do not change. The new arm premise
contains $\theta\bar v=\bar u_i$, as required by the induction hypothesis.
Coverage is preserved: a constructor possible for an assignment in the
new context is possible for its composition in the old one.
For the second assertion induct again. Lookup of $x_j$ uses the $j$th
supplied derivation; other lookups and constants are unchanged. At a
lambda or a case field binder first rename to avoid capture. The supplied
derivations remain valid under extra fresh type and term bindings, and
in a branch under its extra equations, by Lemma~\ref{lem:weakening}.
The induction hypotheses re-establish
the body or arm premises. \textsf{Make}, call and conversion preserve
their type premises and are rebuilt on the substituted children.
\end{proof}

\begin{theorem}[Progress and preservation for the indexed core]\label{thm:safety}
If $\varnothing;\varnothing;\varnothing\vdash_0 e:t$ with $t$ closed,
then $e$ is a value or has a step. If $e\longmapsto e'$ then
$\varnothing;\varnothing;\varnothing\vdash_0 e':t$.
In particular, a ground instance of the lambda in Theorem~\ref{thm:elim}
applied to a closed core term of its argument type has these properties.
There is no assertion here for recursion or arbitrary compiler elaborations.
\end{theorem}
\begin{proof}
At closed types with empty equations, \textsf{Conv} identifies only
identical ground trees. Inspecting the last non-conversion value rule
therefore gives canonical forms: a value at $\bar t\arr u$ is a lambda
of that arity, and a value at $D\,\bar w$ is $K_i(\bar v)$ for some
signature instance $\eta$ with $\eta\bar u_i=\bar w$ and fields typed
at $\eta\bar t_i$. Constants have only base types, disjoint from these
two shapes. For progress induct on typing. A closed variable is
impossible. In any strict argument position the first nonvalue steps by
induction. With all children values a construction is a value; a call
uses the first canonical form and beta-reduces; a case uses the second
and coverage to find the tag. No other expression is left.

For preservation, a step in an evaluation position preserves that child's
type by induction; put it back in the same rule, including the original
arm premises and coverage for a changed scrutinee. For beta-reduction
use inversion for the lambda and the term part of
Lemma~\ref{lem:substitution}. For the case reduction take the
\emph{constructor's} ground instantiation $\eta$ from the value
derivation, not a new arbitrary instantiation from the arm. It satisfies
$\bar w=\eta\bar u_i$, so the type part of Lemma~\ref{lem:substitution}
specializes (Case) to a derivation under $\bar x:\eta\bar t_i$ with
no residual index assumption. Its result is still $t$ because the opened
variables were fresh for $t$. The value fields have exactly these types;
the term part of the lemma gives $e_i[\bar v/\bar x]:t$. Any last
conversion can be reapplied. This is the point at which exporting an
existential witness as the caller's arbitrary type would break the proof.
\end{proof}

\paragraph{What would invalidate a stronger conclusion?}
A non-exhaustive match expression can lack a branch for a well-typed
constructor; an unchecked foreign result need not satisfy the canonical
forms lemma. Example~\ref{ex:fd-height} shows how search can reach a
growing proof's base, not that every finite proof will be found.
The FD argument is conditional on the evidence relation, not a validation
of arbitrary recursive dictionaries or user method laws. The linking
corollary also needs well-typed arguments and compatible class layouts.
The sharing invariant does not model a mutable store. Operational
preservation for rank-polymorphic dictionary wrappers, nonlocal exits,
representation changes and the optimizer needs a semantics accounting for
those operations and their calling conventions. Thus the result is not
``every accepted Coalton program runs without error.'' Its useful consequences
are more specific: which substitutions are justified, which equalities
belong only to an arm, and exactly where an additional invariant is needed.

\paragraph{Source examples for comparison.}
In \src{tests/test-files/define-type-gadt.txt}, case~2 warns about a
missing constructor in a match expression, case~6 rejects an incompatible
index, and case~20 rejects a refutable parameter. For recursive calls,
rank scopes and nonlocal returns, see \src{tests/gadt-tests.ct}.
These examples illustrate behavior; the proofs do not assume their results.

\sources \src{tests/type-inference-tests.lisp}:
\src{test-matching-preserves-target-variables},
\src{test-hidden-ascription-generality-diagnostic}.

\end{document}
